% ============================================================
%  P2 — Density-Calibrated Neural Implicit Field for
%        Ballistic GLAD: From Scalar Porosity to Voxel
%        Occupancy Prediction
%
%  Template: AIP Publishing (revtex4-1, aip class option)
%  Target journal: AIP Advances
%
%  Converted 2026-07-10 from the plain-article-class source
%  (05_THESIS_PAPER_ASSETS/article_draft/P2_density_calibrated_NIF_GLAD.tex)
%  to AIP's revtex4-1 template, at the user's request, targeting AIP
%  Advances (companion submission to P1, targeting JVST-A, same
%  publisher family). Content, figures, tables, and references are
%  unchanged; only the document class, title-block macros, and back
%  matter were converted. Bibliography style switched to aipnum4-1.bst
%  (numbered, AIP house style); \citep/\citet unchanged (revtex4-1
%  loads natbib internally).
%  Author: Oussama Ben Haj Ali (LPMS / ENIT / UTM)
% ============================================================
\documentclass[aip,amsmath,amssymb,reprint]{revtex4-1}

\usepackage{graphicx}
\usepackage{booktabs}
\usepackage{siunitx}
\usepackage{xcolor}
\usepackage{nicefrac}

\graphicspath{{../figures/publication/}}

\renewcommand{\ang}{\ensuremath{\alpha}} % siunitx already defines \ang for angle units,
                                          % but that form (\ang{85}) is never used in this
                                          % document -- all uses here are the bare symbol
                                          % form, so overriding it is safe.
\newcommand{\iou}{\ensuremath{\mathrm{IoU}}}
\newcommand{\nif}{\textsc{NIF}}
\newcommand{\film}{\textsc{FiLM}}

% ============================================================
\begin{document}

\title{Density-Calibrated Neural Implicit Field for Ballistic
  Helical GLAD: Scalar Porosity Prediction, Label-Quality
  Diagnosis, and the Path to Binary Voxel Occupancy}

% [2026-09-09 -- same fix as P1: Martin removed (no contribution), author order set
%  to student / co-supervisor (Paris) / principal supervisor last (not a cotutelle).
%  See P1_ballistic_GLAD_scope_limits_AIP.tex for the full rationale.]
\author{Oussama Ben Haj Ali}
\email{oussama.benhadjali@fst.utm.tn}
\affiliation{Laboratoire de Photovolta\"{i}que et Mat\'{e}riaux
  Semi-conducteurs (LPMS), \'{E}cole Nationale d'Ing\'{e}nieurs de Tunis (ENIT),
  Universit\'{e} de Tunis El Manar (UTM), Tunis, Tunisia}

\author{Bruno Gallas}
\affiliation{Institut des NanoSciences de Paris (INSP), CNRS,
  Sorbonne Universit\'{e}, Paris, France}

\author{Ferid Chaffar Akkari}
\affiliation{Laboratoire de Photovolta\"{i}que et Mat\'{e}riaux
  Semi-conducteurs (LPMS), \'{E}cole Nationale d'Ing\'{e}nieurs de Tunis (ENIT),
  Universit\'{e} de Tunis El Manar (UTM), Tunis, Tunisia}

\date{\today}

% ============================================================
% ---- LONG ABSTRACT (preprint/arXiv version; kept for reference, not used in the
%      submitted manuscript, which requires a short single-paragraph abstract) ----
% Running a full ballistic GLAD simulation for every candidate
% deposition angle is slow, which limits how quickly an
% experimentalist can explore the process window (angle, pitch,
% height) before committing to a deposition run.
% A fast learned surrogate that predicts the resulting film
% morphology directly from the deposition angle $\ang$ -- without
% re-running the simulator -- would let that exploration happen in
% seconds rather than in simulation time, and, if it can also
% resolve where within the film mass is located, would give access
% to voxel-level occupancy maps relevant to connectivity- and
% percolation-sensitive film properties.
% Neural implicit fields (\nif{}) have emerged as a compact,
% coordinate-based representation for continuous 3D density fields
% and are a natural candidate for this role: they are compact,
% differentiable, and can in principle encode arbitrary 3D
% structures.
% We investigate their applicability to ballistic helical
% glancing-angle deposition (GLAD) simulations, where the goal
% is to learn the mapping from deposition angle $\ang$ and spatial
% coordinates $(x,y,z)$ to voxel occupancy.
% We report three successive results.
% First, a separately-configured \film{}-conditioned \nif{} variant
% (``\nif{}~v5''; wider network, an additional held-fixed input
% coordinate, Section~\ref{ssec:exp1}) trained and evaluated on a
% different, older 7-angle dataset restricted to
% $\ang \in \{75, 77.5, 80, 82.5, 85, 86, 87\}^{\circ}$ reconstructs the
% full 3D voxel density field with a mean spatial Pearson
% $r = 0.984$ between predicted and ground-truth fields, averaged
% across these seven angles; this shows the architecture family can
% reconstruct the depth/lateral density profile within this trained
% angular window, and does not establish generalisation to
% $\ang \notin [75^{\circ},87^{\circ}]$ (in particular $60$--$70^{\circ}$ and $89^{\circ}$,
% which this experiment neither trains nor evaluates on).
% Second, the same architecture trained on binary cross-entropy (BCE)
% with Gaussian-smeared occupancy labels yields per-$\alpha$
% intersection-over-union (IoU) indistinguishable from the trivial
% predictor (all alphas: $\Delta\iou < 0.003$), and BCE loss is
% frozen at 0.541 throughout training.
% An $(\alpha,z)$-oracle diagnostic — the theoretically
% optimal depth-aware predictor built directly from training data,
% independently reproduced from the archived dataset (see
% Section~\ref{ssec:res_oracle}) — adds $\delta\iou \leq 0.0075$ over
% the alpha-mean baseline for all main alphas, showing that architecture
% capacity and depth-signal absence alone are unlikely to be the sole
% explanations.
% The root cause is identified as the \emph{labelling scheme}:
% max-normalised bead-count histograms place $90$--$98\%$ of labels
% in the soft zone $(0.1, 0.9)$, eliminating binary boundary information
% from the BCE gradient.
% Third, retraining with hard binary labels (thresholded at $0.5$)
% confirms that the labelling scheme was the bottleneck: BCE decreases
% from $0.669$ (epoch~1) to $0.592$ (best at epoch~50), and the model
% produces genuine spatial predictions with per-$\alpha$ spread of
% $0.06$--$0.13$ in sigmoid output.
% However, only $1/7$ deposition angles ($\ang = 89^{\circ}$, sparse film,
% $\Delta\iou = +0.029$) achieves non-trivial spatial discrimination;
% denser alphas remain trivial predictors at optimal threshold.
% We identify two confounds: multi-seed label conflicts at
% $\ang = 80^{\circ}$ and $85^{\circ}$ (same spatial coordinate with different
% binary labels across seeds) and balanced-sampling distortion of
% the test class distribution.
% These results establish that hard binary labels are a \emph{necessary
% but not sufficient} condition for spatial voxel occupancy prediction.
% A follow-up run combining seed-isolated splits with focal loss
% (early-stopped at epoch~291) was subsequently carried out and did
% \emph{not} solve the spatial-learning problem: aggregate IoU looked
% positive (0.419) but per-$\alpha$ inspection showed the model
% remained statistically indistinguishable from the trivial predictor
% at 5/7 angles ($\Delta\iou < 0.0001$ at 65--85^{\circ}), with only modest,
% non-trivial improvement at $\ang = 87^{\circ}$ ($\Delta\iou = 0.005$) and
% $\ang = 89^{\circ}$ ($\Delta\iou = 0.058$). Three further remedies were
% then tested directly: a multi-resolution hash-grid spatial encoder
% and a continuous-density regression reframing both give no
% meaningful improvement over the hard-label baseline at any angle;
% a partial-realization conditioning approach, which predicts a
% realization's upper film region from an encoded latent of that
% same realization's own near-substrate region, fails a 3-seed
% causal-substitution permutation test at both multi-seed angles, for
% two mechanistically distinct reasons -- an observed region that
% is itself seed-invariant (saturated) at $\ang=80^{\circ}$, and a real but
% under-exploited conditioning signal (encoder influence
% $\sim 10^{-3}$) at $\ang=85^{\circ}$. A direct follow-up retrained this
% conditioning pathway fully unclamped (a 10$\times$ larger allowed
% influence on the output) to test whether the damping itself was the
% cause: it was not -- the permutation test still fails at both angles
% with the same order-of-magnitude same/cross difference, ruling out
% conditioning-pathway capacity as the bottleneck and narrowing it,
% provisionally, to the number of realizations sharing each angle
% (three). We then tested that data-scarcity hypothesis directly: five
% additional $\ang=85^{\circ}$ simulations brought this angle to $n=8$
% realizations, and both variants were retrained from scratch. The
% permutation test still fails at $n=8$ for both variants
% ($p=0.35$ damped, $p=0.53$ unclamped, evaluated over $56$
% cross-realization pairs versus $6$ previously), with same- and
% cross-realization correlation means remaining statistically
% indistinguishable rather than merely falling short of significance.
% More than doubling the per-angle realizations does not resolve Fix
% C's failure, ruling out the practically-achievable data-scarcity
% explanation and leaving the $\ang=85^{\circ}$ result as a genuine, tested
% negative finding.

\begin{abstract}
A neural implicit field (\nif{}) is a compact machine-learning model that
takes a small set of numbers as input -- here, the deposition angle $\ang$
and a point's spatial coordinates $(x,y,z)$ -- and outputs a prediction,
offering a potentially much faster alternative to a full physics simulation
for predicting film structure in ballistic helical glancing-angle
deposition (GLAD). We report three successive attempts to build such a
model, each diagnosing why the previous one fell short. First, a variant
conditioned on a reference simulation via feature-wise linear modulation
(\film{}) successfully reconstructs the full 3D density field, with mean
spatial Pearson correlation $r=0.984$ -- but only for deposition angles
within the $75$--$87^{\circ}$ range it was trained on, not beyond it. This
figure is also in-sample: evaluated at three angles inside that same
range but withheld from training ($78.75^{\circ}$, $83.5^{\circ}$,
$84.25^{\circ}$), the identical checkpoint collapses to $r \approx 0.003$
mean, showing the network reproduces the seven trained density fields
individually rather than learning a function that interpolates smoothly
between them, despite an explicit angle-smoothness regulariser applied
during training. An independent replication of the identical recipe on a
disjoint dataset spanning the full $60$--$89^{\circ}$ deposition-angle
range (7 trained / 2 held-out angles) confirms this is not an artefact of
the narrow trained window: in-sample correlation reaches $r=0.954$ but
collapses to $r=0.088$ at the two held-out angles. Two follow-up tests
narrow down the cause: replacing the \film{} conditioning mechanism
with direct Fourier-encoded concatenation of $\ang$, keeping the same
full-field target, does not help ($r=0.101$ held-out, statistically
indistinguishable from \film{}'s $r=0.088$), while keeping \film{} but
reformulating the \emph{target} as a compact pair-correlation
descriptor $h_g(r;\ang)$ -- the same descriptor family already
validated for P1 -- generalises well on this same wide angular range
(mean $R^2=0.962$, Pearson $r=0.997$ across two held-out angles),
though not clearly better than simple interpolation on the same
points. Together these indicate the failure lives in the
high-dimensional field target, not the conditioning architecture. The
same reframing does not help uniformly: a scalar porosity($\ang$)
target is better served by classical linear interpolation than by any
network at this data volume, and a depth-density-profile descriptor's
generalisation is angle-dependent: after fixing a random-seed
confound in the shared training code (which had been swinging the
headline result between clearly negative and clearly positive across
reruns) and averaging over ten properly-seeded repeats, the profile's
\emph{shape} is reliably recovered ($r=0.766\pm0.041$) but its
absolute level only generalises well inside the training range
($R^2=0.457\pm0.236$ at $72^{\circ}$) and fails near its edge
($R^2=-0.419\pm0.385$ at $87^{\circ}$).
Second, when instead trained to predict binary voxel occupancy (whether a
given point is solid or empty) using softened training labels, the same
architecture collapses to the trivial prediction (always guessing the same
value) at every angle, with essentially no improvement over that trivial
baseline ($\Delta\iou<0.003$). An independently reproduced diagnostic test
-- an $(\alpha,z)$-oracle check -- isolates the cause precisely: this is a
failure of how the training labels were constructed, not a limitation of
the model architecture or a lack of depth information in the input --
$90$--$98\%$ of the softened labels fall in an ambiguous middle range
$(0.1,0.9)$, which removes the sharp boundary signal the model needs to
learn from. Switching to sharp (hard binary) labels restores genuine
spatial learning, but only for the most sparsely-packed film in the dataset
($\ang=89^{\circ}$, improvement $\Delta\iou=+0.029$); denser films remain
unlearnable. We then tested three further fixes -- a higher-resolution
multi-resolution hash-grid input encoding, predicting a continuous density
instead of binary occupancy, and conditioning the model on a partial
(incomplete) realization of the film (damped and fully unclamped \film{}
variants) -- and none produced meaningful spatial learning at the angles
where multiple independent simulation runs were available for testing. For
the partial-realization approach specifically, we checked directly whether
the problem was simply too few independent runs to learn from (three per
angle): enlarging one angle's dataset to eight independent runs and
retraining did not help -- a statistical permutation test still shows no
learned signal for either \film{} variant tested ($p=0.35$ and $p=0.53$,
across $56$ pairs of runs compared), ruling out limited data as the
explanation. This diagnosis is further corroborated by four additional,
architecture-independent checks -- a classical stochastic-reconstruction
baseline, an independently-computed global-connectivity descriptor, a
spatially-aware convolutional network, and a resolution-preserving
local-conditioning variant -- which all converge on the same conclusion:
the missing information is not resident in any low-order structural
descriptor of these films, for any method tested. The result is both a
verified, step-by-step diagnostic
framework for this class of failure in GLAD voxel-occupancy prediction, and
a tested, strengthened negative conclusion: none of the remedies explored
here, including simply providing more training data, resolves the
underlying non-identifiability -- multiple film realizations grown under
identical conditions are genuinely indistinguishable from coordinates and
angle alone, which is exactly the signal a spatial coordinate-based model
would need in order to learn a meaningful occupancy boundary.
\end{abstract}

\maketitle

\medskip
\noindent\textbf{Keywords:} glancing-angle deposition; neural implicit
field; FiLM conditioning; voxel occupancy; label quality; binary
cross-entropy; ballistic simulation; porosity prediction

% ============================================================
\section{Introduction}
\label{sec:intro}

Glancing-angle deposition (GLAD) produces self-assembled nanostructures
whose morphology is exquisitely sensitive to the deposition angle $\ang$,
an angle-controlled columnar growth process reviewed extensively
\citep{hawkeye2014,barranco2016,hawkeye2007}, with self-assembled column
control tracing back to the chiral sculptured thin films of
\citet{robbie1996nature} and \citet{brett2008science}, and treated at
length in \citet{lakhtakia2005stf}.
This sensitivity is exactly what makes GLAD attractive to
experimentalists building functional thin films: within the same
laboratory, angle-tilted Cu$_2$O and CuO$_x$ GLAD films are used as
anisotropic optoelectronic elements \citep{chaffar2018,chaffar2019}
and, more recently, as room-temperature ethanol/methanol gas sensors
whose response depends on the deposition-angle-controlled porosity
and column morphology \citep{benamor2026}.
Literature surveys of GLAD-deposited thin films show a diverse palette
of materials and deposition conditions: Cu and its oxides (Cu$_2$O, CuO)
dominate at high angles ($\ang = 70\text{--}85^{\circ}$), while lower-angle
($\ang < 60^{\circ}$) applications favour W, Ni, and metal oxides (WO$_3$, TiO$_2$),
with temperatures ranging from room-temperature growth to substrate heating
at 200--700\,K across different applications.
Mapping out which angle (and, eventually, which pitch and height)
gives a desired morphology currently requires either a full
experimental deposition run or a full ballistic simulation for every
candidate condition, both of which are slow relative to the size of
the parameter space an experimentalist would like to screen.
Predicting the 3D density field from deposition parameters would
instead enable fast, closed-loop film-design without running a full
ballistic simulation for each condition.
Beyond the scalar porosity trend, knowing \emph{where} within the
film mass is located voxel-by-voxel matters physically: it is a
prerequisite for assessing connectivity and percolation-relevant
properties (e.g.\ whether neighbouring columns touch and form a
continuous conductive/gas-diffusion path, as opposed to remaining
isolated), which is precisely the kind of structural information
that governs the electrical and sensing behaviour reported in
\citep{chaffar2019,benamor2026}.

\textbf{Relationship to a companion manuscript.} A companion paper
\citep{companionP1} characterises the same ballistic simulator's void fraction,
morphology, and pair-correlation behaviour across a $60$--$89^{\circ}$
helical GLAD campaign, and includes, as one part of that broader
characterisation, a single memorization-style sanity check on
spatial voxel learnability: a coordinate-conditioned NIF trained
and evaluated on one deterministic realization achieves
near-perfect IoU, which that paper uses to argue that the observed
production-training BCE plateau is not, by itself, evidence of
representational impossibility.
The present paper is a dedicated, considerably deeper investigation
of specifically the voxel-occupancy learnability question raised by
that plateau: rather than a single overfit test, it systematically
diagnoses \emph{why} naive spatial learning fails (soft-label
information loss, Section~\ref{ssec:disc_softtrap}) and tests, and
rules out, several standard remedies (mass-conservation penalty,
hard labels, focal loss, seed isolation).
The two studies are complementary, not duplicative: the companion
paper's single deterministic overfit result is, in fact, a data
point that this paper's own diagnostic framework predicts should be
achievable (a clean, single-realization, non-pooled label is exactly
the condition under which the soft-label equilibrium trap identified
here does not apply — see the cross-reference to that result in
Section~\ref{ssec:disc_softtrap}); it is not a contradiction of this
paper's still-largely-negative overall result on the harder,
pooled multi-seed, multi-angle production setting.
We do not claim this paper closes the voxel-occupancy question — the
Conclusion (Section~\ref{sec:conclusion}) is explicit that a higher-capacity
local spatial encoding, once tested here, also fails to resolve it — only
that each successive test narrows the failure
to a specific, mechanistically understood cause.

Neural implicit fields, which represent a continuous scalar field as
a function of coordinate inputs \citep{mildenhall2020nerf}, are a
natural candidate: they are compact, differentiable, and can in
principle encode arbitrary 3D structures.

The \film{} conditioning mechanism \citep{perez2018film} allows a
single network to represent fields parameterised by an external
condition — here, the deposition angle $\ang$.
This paper pursues three objectives, framed around what an
experimentalist would want such a surrogate to eventually deliver:
\begin{enumerate}
  \item \textbf{Scalar porosity surrogate.} Can a \film{}-conditioned
    \nif{} learn the $\ang$-dependent scalar porosity of ballistic
    GLAD helical films — i.e.\ replace a full ballistic run with a
    fast lookup for the single number (film openness/density) that
    experimentalists typically track first?
  \item \textbf{Voxel-level occupancy.} Can it go further and predict
    binary voxel occupancy (occupied vs.\ empty) at given spatial
    coordinates — the finer-grained information needed for
    connectivity/percolation-relevant structural analysis, not just
    the bulk average?
  \item \textbf{Root-cause diagnosis.} If voxel-level prediction
    fails, what is the root cause, and can it be remedied?
\end{enumerate}
The Conclusion (Section~\ref{sec:conclusion}) answers these three
objectives in the same order.

We answer these questions through three sequential experiments.
Experiments~2 and~3 (spatial voxel occupancy) share a common
dataset of twelve ballistic Cu helical GLAD configurations
($h = \SI{300}{nm}$, pitch $= \SI{150}{nm}$,
$\ang \in \{60, 65, 70, 75, 80, 85, 87, 89\}^{\circ}$).
Experiment~1 (scalar porosity) predates this dataset and uses a
separate, older 7-angle campaign restricted to
$\ang \in \{75, 77.5, 80, 82.5, 85, 86, 87\}^{\circ}$
(Section~\ref{ssec:exp1}); its result therefore speaks only to
that narrower angular window, not to the full
$\{60,\ldots,89\}^{\circ}$ range used in Experiments~2--3.

% ============================================================
\section{Methods}
\label{sec:methods}

\subsection{Dataset and voxel representation}
\label{ssec:data}

Every surrogate trained below is only as good as the simulation
data it is trained on, so before describing any network we first
describe how the ballistic simulator's raw output (a cloud of
deposited bead positions) is turned into the regular 3D grid of
occupancy values the network actually learns from.
Simulations are performed with GLAD~v3.2 (bead-sphere ballistic,
no surface diffusion) on a $100 \times 100\,\mathrm{nm}$ periodic
substrate, producing 12 configurations.
Atom positions are binned onto a
$50 \times 50 \times 151$ voxel grid ($\SI{2}{nm}$ isotropic
resolution) and normalised by the maximum bin count,
yielding a density field $\rho_{\rm vox}(\mathbf{r}) \in [0,1]$.
This max-normalised bead-count histogram encodes the local
occupation density relative to the densest voxel in the film.
Density and occupancy, rather than column-fanning or pitch descriptors,
are used as the surrogate targets throughout this paper because a
dedicated single-condition ablation ($\alpha=85^{\circ}$, seed~0,
$300\to\SI{700}{nm}$, eight height stages) found that switching on a
basic, uncalibrated-mobility surface-diffusion step changes the atom
density ratio by only $0.9995$--$1.0021$ relative to the diffusion-off
baseline and leaves the fanning, helical-pitch, and column-centreline
descriptors statistically indistinguishable from the diffusion-off case
at every height (the same watershed descriptor separately recovers
$p\approx0.28$--$0.34$ on a synthetic separated-column control, so the
null result reflects the simulator's fused morphology rather than a
broken descriptor). Within this reduced-order model, fanning is
therefore not a robust learning target at the angles and heights used
here, while density and occupancy are well-defined and change smoothly
with $\ang$; this motivates the choice of surrogate target below and
is not a claim that surface diffusion is irrelevant to copper GLAD in
general.
A follow-up extension used exclusively by the partial-realization-conditioning test
(Section~\ref{ssec:disc_tier2}) added five further realizations at $\ang=85^{\circ}$
(seeds 3--7, identical protocol), bringing that angle to $n=8$; all other experiments
use the twelve-configuration dataset above.
% [2026-09-08, per O.B.H.A. instruction -- scope-freeze provenance note; see P2_SCOPE_FREEZE_20260908.md]
All voxel fields used here were generated with the simulator core as it stood when this
dataset was archived. A subsequent grid-rebuild correction to that core (documented in
the companion paper) post-dates this dataset and changes per-angle atom counts by up to a
few percent, without altering the invariant topology class. Because every result in this
paper is a statement about \emph{independent realizations of the same growth condition}
being mutually indistinguishable from coordinates and angle alone -- a property of the
ensemble, not of the specific core version -- regenerating the dataset on the corrected
core is not expected to change any conclusion; it is left as future work, and the numbers
below should be read as pertaining to the archived dataset.

\subsection{FiLM-conditioned NIF architecture (Experiments 2--3)}
\label{ssec:arch}

The goal of this architecture is to let a single trained network
answer ``is this point occupied?'' at any spatial coordinate for
any deposition angle in the training range, i.e.\ to act as a
continuous, queryable replacement for a full-grid ballistic
simulation output.
This subsection describes the architecture shared by
Experiments~2 and~3.  Experiment~1 uses a separately configured
network variant, described in its own subsection
(Section~\ref{ssec:exp1}), which is \emph{not} identical to the
one below (different hidden width, Fourier-frequency count, and
input coordinates).
The network receives a coordinate vector
$\mathbf{x} = (x_{\rm n}, y_{\rm n}, z_{\rm n}, \alpha_{\rm n}) \in \mathbb{R}^4$,
where spatial coordinates are unit-normalised to $[0,1]$ and
$\alpha_{\rm n} = (\ang - 74.5)/14.5$ maps the deposition-angle
range $[60^{\circ}, 89^{\circ}]$ to approximately $[-1, 1]$.
This coordinate-to-scalar architecture follows the broader family of
neural implicit fields established for shape and scene representation
\citep{park2019deepsdf,mescheder2019occupancy,mildenhall2020nerf}; a
periodic-activation alternative to the Fourier-encoding-plus-ReLU
approach used here is given by \citet{sitzmann2020siren}, and hash-grid
spatial encodings \citep{muller2022instant} were tested as an alternative
input representation for this simulator (Section~\ref{ssec:disc_tier2}).
A Fourier positional encoding \citep{tancik2020fourier} with
$n_{\rm freq} = 32$ frequencies encodes $(x_{\rm n}, y_{\rm n}, z_{\rm n})$;
such coordinate-based MLPs are known to be biased toward learning
low-frequency function components first (the ``spectral bias'' of
\citet{rahaman2019spectral}), which the Fourier encoding is specifically
intended to counteract.
The encoded vector is passed through 4 hidden layers of width 256
with \film{} conditioning at each layer on the $\alpha_{\rm n}$
scalar, implemented as element-wise affine scaling $\gamma$ and
shift $\beta$ predicted by a two-layer sub-network.
The final output is a sigmoid-activated scalar $\hat{\rho} \in [0,1]$.
Figure~\ref{fig:nif_arch_p2} summarises this shared architecture.

\begin{figure}[htbp]
  \centering
  \includegraphics[width=\linewidth]{nif_architecture_p2}
  \caption{%
    Experiments 2 and 3 differ only in what target this same network is
    trained against (continuous soft labels versus binarised hard
    labels); the architecture and the FiLM conditioning mechanism itself
    are unchanged between them, which is why the comparison between the
    two experiments isolates the effect of the labelling scheme.
    Schematic of the shared \film{}-conditioned NIF architecture used
    for Experiments 2--3 (conceptual diagram; Experiment 1 uses a
    separately configured variant, Section~\ref{ssec:exp1}).
    A spatial coordinate is Fourier-encoded and passed through four
    hidden layers, each modulated by a deposition-angle-conditioned
    scale $\gamma_i$ and shift $\beta_i$ (\film{}); the output is a
    sigmoid-activated occupancy/density prediction $\hat{\rho}$.
  }
  \label{fig:nif_arch_p2}
\end{figure}

\subsection{Experiment 1: full-field density reconstruction (NIF~v5)}
\label{ssec:exp1}

This experiment asks whether a fast surrogate can replace an
expensive ballistic run for reproducing the \emph{continuous}
density field itself (the depth/lateral density profile an
experimentalist would compare to a cross-sectional density or
porosity measurement), rather than a single scalar summary of it.
This experiment predates Experiments~2--3 and uses a distinct,
separately-frozen network configuration (script
\texttt{train\_nif\_v5\_film\_smooth.py}), not the shared
architecture of Section~\ref{ssec:arch}: a 4-D coordinate input
$(x_{\rm n}, y_{\rm n}, z_{\rm n}, t_{\rm n})$ with $z$ normalised
against $Z_{\max} = \SI{700}{nm}$ (not \SI{300}{nm}) and a fourth
coordinate $t$ held fixed at its maximum value for every sample in
this configuration (a vestigial input carrying no information
here), a Fourier positional encoding \citep{tancik2020fourier} with
64 frequencies per scale (vs.\ 32 in Section~\ref{ssec:arch}), and
hidden width 512 (vs.\ 256);
deposition angle is injected via the same \film{} mechanism, using
a disjoint 7-angle training set
$\ang \in \{75, 77.5, 80, 82.5, 85, 86, 87\}^{\circ}$ (3 angles —
$78.75^{\circ}, 83.5^{\circ}, 84.25^{\circ}$ — held out).
The network is trained to reconstruct the full 3D voxel density
field $\rho_{\rm vox}(x,y,z)$ for each training angle (not a
1D scalar-per-angle regression target).
Loss: mean-squared error over voxels.
Metric: for each of the 7 training angles, the spatial Pearson
correlation between the predicted and ground-truth full 3D density
field is computed, and the reported $r$ is the mean of these 7
per-angle spatial correlations — a measure of how well the model
reconstructs the depth/lateral density profile at each trained
angle, not a cross-angle scalar-trend regression accuracy.

\subsection{Experiment 2: soft-label spatial occupancy (PINN-BCE)}
\label{ssec:exp2}

This experiment tests whether the same network can go beyond a
continuous density value and predict, at each spatial coordinate,
whether that point is physically occupied by material or not — the
voxel-level occupancy information that connectivity/percolation
analysis needs — using the max-normalised density value directly as
a soft occupancy target, and adding a physics-motivated penalty
that keeps the predicted mean occupancy consistent with the
simulator's measured mass per angle.
Figure~\ref{fig:vox_concept} illustrates, on a small synthetic example,
why this task is trivial in one regime and impossible in another: a
classifier can learn a sharp occupied/empty boundary, but has nothing
to fit if almost every voxel carries an intermediate, ambiguous value.
Section~\ref{ssec:res_oracle} shows that the real dataset for this
simulator sits much closer to the second case at most angles.

\begin{figure}[htbp]
  \centering
  \includegraphics[width=0.85\linewidth]{voxel_occupancy_concept}
  \caption{%
    The distinction driving this experiment's outcome in one picture:
    voxel-occupancy prediction is easy exactly when the labels are
    already almost binary, and structurally impossible when they are
    not, regardless of network capacity.
    Conceptual illustration (synthetic $10\times10$ grid, not
    simulation data) of the voxel-occupancy prediction task.
    (a) A sharp occupied/empty boundary: a binary classifier has a clear
    decision boundary to learn.
    (b) An ambiguous field where most voxels sit near the intermediate
    value $0.5$ rather than at either extreme: there is no clean
    boundary for a classifier to recover, independent of architecture.
  }
  \label{fig:vox_concept}
\end{figure}
The network is trained on 600\,000 random voxel samples
($n_{\rm train} = 480\,000$, $n_{\rm val} = 60\,000$,
$n_{\rm test} = 60\,000$) from the max-normalised density fields.
Each sample is a tuple $(x_{\rm n}, y_{\rm n}, z_{\rm n},
\alpha_{\rm n}, \rho_{\rm vox})$ where $\rho_{\rm vox}$ serves
as the soft BCE target.
A mass-conservation penalty $\lambda_{\rm mass} = 10.0$, in the spirit of
the physics-informed residual penalties of \citet{raissi2019pinn}, is
added to align predicted mean occupancy per $\ang$ with the measured
ground-truth occupied fractions.
Loss: $\mathcal{L} = \mathrm{BCE} + \lambda_{\rm mass}\,\mathcal{L}_{\rm mass}$.
Best-model criterion: validation IoU at threshold 0.5.
Training is warm-started from a separate 200-epoch BCE-only
pretraining run on the same architecture and dataset
(no mass-conservation term, best validation IoU $=0.183$); the
mass-conservation fine-tuning stage reported here raises this to
$\iou = 0.242$ (Section~\ref{ssec:res2}), a $+32\%$ relative
improvement attributable to the mass-conservation term.

\subsection{Experiment 3: hard-label spatial occupancy}
\label{ssec:exp3}

Following the diagnosis in Section~\ref{ssec:oracle} that the
continuous, max-normalised occupancy target itself was the
obstacle (Section~\ref{ssec:res_oracle}), this experiment tests
whether giving the network an unambiguous occupied/empty label at
each voxel is enough to recover genuine spatial (connectivity-
relevant) discrimination, rather than a value diluted across a
Gaussian-smeared range.
The density field is binarised:
$y_{\rm hard}(\mathbf{r}) = \mathbb{1}[\rho_{\rm vox}(\mathbf{r}) > 0.5]$.
The dataset is rebuilt from the same voxel files using identical
sampling parameters (50\,000 samples per configuration, balanced
$\nicefrac{1}{3}$ occupied / $\nicefrac{1}{3}$ empty /
$\nicefrac{1}{3}$ random) with hard binary targets.
The \film{}-conditioned \nif{} is initialised from scratch (no
warm-start from the soft-label checkpoint, which had converged to
predicting $\hat{\rho} \approx 0.5$ everywhere).
Training uses pure BCE on binary $\{0,1\}$ targets.
Best-model criterion: validation IoU at threshold 0.5.

\subsection{Oracle predictor and root-cause diagnostic}
\label{ssec:oracle}

Before concluding that the network itself is incapable of spatial
prediction, it is important to rule out a simpler explanation: that
depth alone (independent of any learned lateral $(x,y)$ structure)
already carries most of the achievable signal, in which case no
architecture change would help.
To distinguish architecture failure from label-quality failure, we
construct an $(\alpha, z)$-oracle predictor: for each test voxel,
the prediction is the mean training $\rho_{\rm vox}$ in the same
$(\ang, z_{\rm bin})$ cell, where $z_{\rm bin}$ resolves the
depth at $\SI{2}{nm}$ resolution.
This oracle captures the maximum information extractable from the
depth profile alone, without any spatial $(x,y)$ discrimination.
IoU, the standard overlap metric for binary spatial segmentation tasks
\citep{rezatofighi2019giou}, is evaluated after a threshold sweep on the
oracle predictions;
a $\delta\iou \leq 0.01$ gap relative to the alpha-mean trivial
baseline is taken as evidence that depth carries no actionable
binary boundary signal in the soft labels.

% ============================================================
\section{Results}
\label{sec:results}

\subsection{Experiment 1 — full-field reconstruction at trained angles: strong positive result, but no angular interpolation}
\label{ssec:res1}

This subsection addresses Objective~1 (scalar/continuous-field
porosity surrogate; Section~\ref{sec:intro}).
\nif{}~v5 achieves a mean spatial Pearson $r = 0.984$ between
predicted and ground-truth full 3D density fields, averaged over
the 7 trained angles $\ang \in \{75, 77.5, 80, 82.5, 85, 86, 87\}^{\circ}$
(Section~\ref{ssec:exp1}), confirming that this \film{}-conditioned
variant successfully reconstructs the depth/lateral density
profile within this angular window.
This establishes a positive baseline for that window only: it does
not demonstrate cross-angle scalar-trend regression, and it makes
no claim about $\ang \notin [75^{\circ},87^{\circ}]$, which this experiment does
not train or evaluate on.
Table~\ref{tab:nifv5_peralpha} shows the per-angle breakdown behind
the aggregate: the spatial Pearson correlation is uniformly high
($r \geq 0.980$) across all 7 trained angles, with no systematic
degradation toward either edge of the trained window, and the mean
relative density error stays below $1.3\%$ at every angle.

\begin{table}[htbp]
  \centering
  \caption{%
    The strong aggregate correlation reported above is not driven by
    a few easy angles: every one of the 7 trained angles individually
    reaches $r \geq 0.98$.
    Per-angle full-field reconstruction quality for \nif{}~v5
    (best checkpoint, epoch 1000). $r$: spatial Pearson correlation
    between predicted and ground-truth 3D density fields.
    Mean error: mean absolute relative density error (\%).
  }
  \label{tab:nifv5_peralpha}
  \begin{tabular}{rcc}
    \toprule
    $\ang$ (°) & Pearson $r$ & Mean error (\%) \\
    \midrule
    75.0 & 0.9857 & 0.065 \\
    77.5 & 0.9866 & 0.346 \\
    80.0 & 0.9850 & 0.028 \\
    82.5 & 0.9850 & 0.899 \\
    85.0 & 0.9820 & 0.775 \\
    86.0 & 0.9804 & 0.472 \\
    87.0 & 0.9833 & 1.203 \\
    \midrule
    Mean & 0.9840 & 0.541 \\
    \bottomrule
  \end{tabular}
\end{table}

Figure~\ref{fig:nifv5_slice} makes this agreement directly visible rather
than only numerical: a vertical slice through the predicted and
ground-truth fields at $\ang=80^{\circ}$ shows the same fine tilted-column
texture reproduced by the network, with a residual that is visually flat
noise rather than a structured bias.

\subsubsection{Held-out angular interpolation: confirmed negative result}
\label{sssec:res1_heldout}

All results reported above are evaluated at the same 7 angles used for
training, i.e.\ in-sample. Three additional density fields were
simulated at angles chosen strictly \emph{inside} the trained
$[75^{\circ},87^{\circ}]$ window but never shown to the network during
training -- $\ang=78.75^{\circ}$ (between the $77.5^{\circ}$ and
$80^{\circ}$ training points), $83.5^{\circ}$, and $84.25^{\circ}$
(both between the $82.5^{\circ}$ and $85^{\circ}$ training points) --
to test whether the \film{}-conditioned network has learned a
genuinely angle-continuous density function or has instead memorised
seven discrete fields. We ran the saved best checkpoint
(epoch 1000, the same checkpoint behind Table~\ref{tab:nifv5_peralpha})
forward at these three interpolated angles, using the identical
coordinate normalisation, \film{} conditioning, and spatial Pearson
metric as the in-sample evaluation, so the two numbers are directly
comparable.
Table~\ref{tab:nifv5_heldout} reports the result: mean spatial Pearson
correlation collapses from $r=0.984$ in-sample to $r=0.003$ at the
three held-out angles -- consistent with zero correlation at every
individual angle (all three individual $|r| < 0.007$), i.e.\ no usable
spatial correlation at all -- even though the network's predicted mean
density stays close to the ground-truth mean density at each angle
(relative error $0.5$--$3.6\%$). This combination is diagnostic: the
network gets the bulk (zeroth-moment) density right but has no access
to the correct spatial (pointwise) structure at an untrained angle,
which rules out a merely quantitative degradation and instead
indicates that the seven trained fields are each represented
individually rather than as samples of one smooth underlying function
of $\ang$, despite the smoothness penalty (Section~\ref{ssec:arch})
applied during training. We report this result exactly as obtained,
with no retraining or checkpoint selection after the fact.

\begin{table}[htbp]
  \centering
  \caption{%
    Held-out angular-interpolation test for \nif{}~v5 (identical
    checkpoint as Table~\ref{tab:nifv5_peralpha}, epoch 1000, no
    retraining). The 3 angles were simulated but withheld from
    training; each lies strictly inside the trained
    $[75^{\circ},87^{\circ}]$ window. $r$: spatial Pearson correlation
    between predicted and ground-truth 3D density fields. Mean error:
    mean absolute relative density error (\%).
  }
  \label{tab:nifv5_heldout}
  \begin{tabular}{rcc}
    \toprule
    $\ang$ (°) & Pearson $r$ & Mean error (\%) \\
    \midrule
    78.75 & 0.0067 & 0.519 \\
    83.5  & $-0.0045$ & 2.231 \\
    84.25 & 0.0063 & 3.619 \\
    \midrule
    Mean & 0.0028 & 2.123 \\
    \bottomrule
  \end{tabular}
\end{table}

\begin{figure}[htbp]
  \centering
  \includegraphics[width=\linewidth]{nifv5_pred_vs_truth}
  \caption{%
    The high Pearson correlation in Table~\ref{tab:nifv5_peralpha} is not
    a coincidence of a coarse aggregate statistic: the network reproduces
    the fine, tilted-column texture of the ground-truth field directly,
    and the residual is visually featureless.
    Vertical ($xz$) slice through the trained NIF~v5 field at
    $\ang=80^{\circ}$ (a training angle), generated by a single forward
    pass through the saved checkpoint (no retraining).
    (a) Ground-truth density field.
    (b) NIF~v5 prediction.
    (c) Pointwise difference (predicted $-$ truth); $r=0.985$ for this
    slice, consistent with Table~\ref{tab:nifv5_peralpha}.
  }
  \label{fig:nifv5_slice}
\end{figure}

\subsubsection{Wide-range replication (60--89\textdegree{}): held-out interpolation failure confirmed on an independent dataset}
\label{sssec:res1_widerange}

The held-out interpolation collapse above was measured on the same
narrow $[75^{\circ},87^{\circ}]$ window, $700\,$nm-height,
static-rotation dataset used to train \nif{}~v5, leaving open whether
the failure is specific to that narrow window or a more general
property of the \film{}-conditioned recipe. We re-ran the identical
architecture (multi-scale Fourier features, \film{} $\alpha$-conditioning,
MSE loss, angular-smoothness penalty; Section~\ref{ssec:arch}) on an
independent, already-simulated dataset spanning the \emph{full}
physical range used elsewhere in this project: box $=150\,$nm,
target height $=300\,$nm, pitch $=150\,$nm, helical rotation, seed~000,
at $\ang \in \{60,65,70,72,75,80,85,87,89\}^{\circ}$. Train/held-out
angles were pre-registered before any training or evaluation:
\textbf{train} $\ang \in \{60,65,70,75,80,85,89\}^{\circ}$ (7 angles,
spanning the full range), \textbf{held out}
$\ang \in \{72,87\}^{\circ}$ (2 angles, interior to the trained range,
never loaded until the pre-declared epoch budget was reached).

Training ran for 60 epochs (cosine-annealed learning rate,
$5\times10^{-4}\to10^{-5}$) rather than the 400-epoch budget
originally planned for this dataset: this environment's GPU idles at
$\sim\!72$--$78^{\circ}$C, and sustained training load reaches the
project's $85^{\circ}$C thermal-safety threshold within about a
minute, so training was run duty-cycled (paused above $85^{\circ}$C,
resumed below $80^{\circ}$C, hard-killed above $90^{\circ}$C as a
safety net -- never triggered), which caps realistic throughput to
roughly one epoch per 2.4 minutes on this hardware. The epoch budget
was reduced accordingly and is reported here as a compute-constrained
deviation from the original plan, not a result-driven one: it was
fixed before the held-out alphas were touched.

At this budget the in-sample fit is strong and comparable to
\nif{}~v5's: mean spatial Pearson $r = 0.954$ across the 7 trained
angles (individual range $0.896$--$0.983$; degrading somewhat at the
$89^{\circ}$ edge of the range, mean density error $2.5\%$ overall).
The held-out result, however, reproduces the same qualitative failure
as Table~\ref{tab:nifv5_heldout}: mean spatial Pearson collapses to
$r = 0.088$ at the two held-out angles
(Table~\ref{tab:widerange_heldout}), with mean relative density error
of $38.5\%$ -- an order of magnitude worse than the in-sample error.
By the same diagnostic standard used in
Section~\ref{sssec:res1_heldout} ($r \lesssim 0.1$ indicating no
usable spatial correlation), this is a second, independent
confirmation of the same negative result: a strong in-sample fit
($r>0.95$) coexisting with a near-total loss of spatial correlation
at nearby, untrained angles. The failure is not an artefact of the
narrow $[75^{\circ},87^{\circ}]$ window, the $700\,$nm/static-rotation
convention, or the specific 7 trained angles used in \nif{}~v5: it
reproduces on a different box size, height, rotation convention, and
angular range, using an entirely disjoint train/held-out split. We
report this exactly as obtained, with no retraining or checkpoint
selection after the fact.

\begin{table}[htbp]
  \centering
  \caption{%
    Held-out angular-interpolation test for the wide-range replication
    (best checkpoint, epoch 60, no retraining), on the independent
    box150\,nm/height300\,nm/pitch150\,nm helical dataset. The 2 angles
    were simulated but withheld from training. $r$: spatial Pearson
    correlation between predicted and ground-truth 3D density fields.
    Mean error: mean absolute relative density error (\%).
  }
  \label{tab:widerange_heldout}
  \begin{tabular}{rcc}
    \toprule
    $\ang$ (°) & Pearson $r$ & Mean error (\%) \\
    \midrule
    72.0 & 0.1699 & 29.85 \\
    87.0 & 0.0063 & 47.22 \\
    \midrule
    Mean & 0.0881 & 38.54 \\
    \bottomrule
  \end{tabular}
\end{table}

\subsubsection{Reformulating the prediction target: compact descriptors generalize where the full field does not}
\label{sssec:res1_descriptors}

Sections~\ref{sssec:res1_heldout}--\ref{sssec:res1_widerange} show that
the raw 3D density field is not a learnable function of $\ang$ at
held-out angles, regardless of dataset. Two follow-up experiments test
whether this failure is caused by \emph{how} the model is asked to
represent the field (its architecture), or by \emph{what} it is asked
to represent (the full voxel field itself, as opposed to a compact
summary of it). A third follow-up tests whether this same descriptor
reframing helps for two further quantities of interest, with mixed
results reported honestly below.

\paragraph{Is it the target, or the conditioning mechanism?}
The wide-range \film{}-conditioned network (Section~\ref{sssec:res1_widerange})
was retrained with an architecturally different alpha-conditioning
scheme, keeping the full 3D density field as the prediction target,
the dataset, and the train/held-out split ($\ang\in\{60,65,70,75,80,85,89\}^{\circ}$
trained, $\ang\in\{72,87\}^{\circ}$ held out) unchanged. Instead of
\film{}'s weight modulation (a small $\alpha$-dependent network
producing per-layer gain/bias terms, hypothesised to degenerate into a
per-angle lookup table), $\alpha$ was Fourier-encoded with the same
recipe used for $(x,y,z)$ and simply concatenated with the encoded
spatial coordinates before an ordinary MLP -- removing weight
modulation entirely. If \film{}'s conditioning mechanism were the
cause of the held-out collapse, this substitution should measurably
improve held-out correlation.

It does not. Held-out mean Pearson correlation is
$r=0.101$ for the concatenation variant versus $r=0.088$ for the
\film{} baseline (Table~\ref{tab:concat_heldout}) -- both effectively
zero and not a meaningful difference given only two held-out angles.
At $\ang=87^{\circ}$, the harder of the two held-out angles, the
concatenation variant is actually worse ($r=0.001$ vs.\ $0.006$) and
its density-magnitude error is substantially worse ($83.6\%$ vs.\
$47.2\%$); its in-sample fit is also uniformly weaker than \film{}'s
(mean trained-angle $r=0.869$ vs.\ $0.954$). This rules out the
conditioning \emph{mechanism} as the primary cause of the held-out
failure: removing \film{}'s weight modulation neither fixes
generalisation nor even matches the baseline's in-sample fit. The
difficulty therefore appears to reside in the high-dimensional,
spatially-detailed full-3D-field \emph{target} itself, not in how
$\alpha$ is injected into the network -- directly motivating the
descriptor-reframing test below.

\begin{table}[htbp]
  \centering
  \caption{%
    Held-out comparison of \film{} weight-modulation conditioning
    against direct Fourier-concatenation of $\alpha$, same dataset,
    split, and full-3D-density-field target as
    Table~\ref{tab:widerange_heldout}. $r$: spatial Pearson
    correlation, held-out angles only.
  }
  \label{tab:concat_heldout}
  \begin{tabular}{rcc}
    \toprule
    $\ang$ (°) & \film{} $r$ & Concat $r$ \\
    \midrule
    72.0 & 0.170 & 0.200 \\
    87.0 & 0.006 & 0.001 \\
    \midrule
    Mean & 0.088 & 0.101 \\
    \bottomrule
  \end{tabular}
\end{table}

\paragraph{Compact descriptors of the same dataset.}
A separate follow-up kept the \film{}-conditioned architecture but
changed the prediction target: instead of the raw $(x,y,z)$-indexed
density field, the network predicts one of three compact,
physically-meaningful descriptors as a function of $\ang$ and a
single 1D coordinate, computed from the identical box150\,nm/
height300\,nm dataset and train/held-out split used in
Section~\ref{sssec:res1_widerange}: (i) $h_g(r;\ang)$, the
corotating-frame pair-correlation function -- the same descriptor and
extraction pipeline P1's own \textsc{CorrNIF} network already uses
successfully on a narrower angular window; (ii) $\rho_z(z;\ang)$, the
mean-density depth profile; and (iii) scalar porosity$(\ang)$. All
three are $\geq 1000\times$ more compact than the raw voxel field
($25$, $152$, and $1$ dimensions respectively, versus $\sim 380{,}000$
voxels).

$h_g(r;\ang)$ generalises well at both genuinely held-out angles:
mean $R^2=0.962$, mean Pearson $r=0.997$ across $\ang\in\{72,87\}^{\circ}$
(Table~\ref{tab:descriptor_heldout}). This replicates, on the wider
$60$--$89^{\circ}$ range and an independent dataset, the same
qualitative result P1's \textsc{CorrNIF} already established on its
own narrower window ($R^2=0.999$ at $\ang=72^{\circ}$): statistical
descriptors of the density field are smooth, learnable functions of
$\ang$ even where the exact voxel realisation, and even the full
3D field's absolute density level, are not. This is the clearest
positive finding of this reformulation test. It should, however, be
read with the same caveat that already applies to P1's own result: the
NIF's $h_g(r)$ held-out fit does not clearly outperform naive linear
interpolation between the bracketing trained angles ($R^2=0.972$ mean)
or nearest-neighbour lookup ($R^2=0.989$ mean) at these same two held-out
points. The finding here is therefore best read as a confirmatory
replication of P1's descriptor-based approach in a new, wider-range
setting -- not as evidence that a NIF outperforms simple interpolation
for this particular descriptor. What the result does establish is that
$h_g(r;\ang)$, unlike the raw field, is the kind of smooth quantity for
which \emph{any} reasonable interpolant (learned or classical) works --
exactly the property the full 3D field lacks.

\begin{table}[htbp]
  \centering
  \caption{%
    Held-out generalisation of three compact descriptors of the same
    box150\,nm/height300\,nm dataset and train/held-out split as
    Table~\ref{tab:widerange_heldout}, predicted by a \film{}-conditioned
    network with the descriptor's own coordinate ($r$ for $h_g$, $z$ for
    $\rho_z$; none for scalar porosity) as input instead of $(x,y,z)$.
    Porosity is reported as mean absolute error (a fractional occupancy,
    not a correlation).
  }
  \label{tab:descriptor_heldout}
  \begin{tabular}{lccc}
    \toprule
    Descriptor & $\ang=72^{\circ}$ & $\ang=87^{\circ}$ & Mean \\
    \midrule
    $h_g(r;\ang)$, $R^2$ & 0.933 & 0.991 & \textbf{0.962} \\
    $h_g(r;\ang)$, Pearson $r$ & 0.999 & 0.996 & \textbf{0.997} \\
    \addlinespace
    $\rho_z(z;\ang)$, $R^2$ & 0.444 & $-0.964$ & $-0.260$ \\
    $\rho_z(z;\ang)$, Pearson $r$ & 0.742 & 0.692 & 0.717 \\
    \addlinespace
    porosity$(\ang)$, mean abs.\ err. & 0.254 & 0.108 & \textbf{0.181} \\
    \bottomrule
  \end{tabular}
\end{table}

\paragraph{Scalar porosity($\ang$): the descriptor reframing does not help here -- classical interpolation wins.}
The scalar porosity$(\ang)$ NIF fails outright (mean absolute error
$0.181$, and predicts a physically impossible value $>1$ at
$\ang=72^{\circ}$), while simple linear interpolation between the two
bracketing trained angles achieves mean absolute error $0.035$ and
cubic-spline interpolation $0.040$ (a follow-up check, same
held-out pair). With only 7 training points and a
non-monotonic porosity($\ang$) curve, this is not evidence against the
descriptor-reframing idea in general -- it is evidence that a learned
network is the wrong tool for this specific quantity at this data
volume. For scalar porosity specifically, classical interpolation is
not merely adequate but strictly better than the learned alternative,
and is the recommended approach for any downstream use of this
quantity.

\paragraph{Depth-density profile $\rho_z(z;\ang)$: shape generalises, absolute scale does not, with a real angle-dependent asymmetry.}
A random-seed confound in this descriptor-reframing training script
was identified and fixed: linear-layer weight initialisation was
drawing from an unseeded generator, which had been swinging this
particular result between clearly negative and clearly positive
across otherwise-identical reruns. All linear-layer weight
initialisation was fixed to draw from the same seeded generator as
the rest of the recipe (verified deterministic: repeat runs under an
identical seed now reproduce bit-identical loss curves). The
held-out evaluation was repeated over ten independent, properly-seeded
runs (seeds 101--110) on the same pre-registered $\{72,87\}^{\circ}$
held-out split. The result is now stable and reproducible, and reveals
a more specific picture than either a clean positive or negative: mean
Pearson $r=0.766\pm0.041$ across all ten seeds (the profile
\emph{shape} is reliably learned, low variance, positive in every
run), while mean $R^2=0.019\pm0.191$ genuinely spans zero (4 of 10
seeds positive, 6 negative) -- the absolute density level is not
reliably recovered. Splitting by held-out angle shows this is not
uniform noise but a real, reproducible asymmetry: at
$\ang=72^{\circ}$ (well inside the trained range) $R^2=0.457\pm0.236$,
mostly positive; at $\ang=87^{\circ}$ (the edge of the trained range)
$R^2=-0.419\pm0.385$, mostly negative. We report this as the
paper's genuine, final finding for this descriptor: a partial success
--- generalisation degrades specifically near the boundary of the
training range, consistent with an extrapolation-proximity effect
rather than an unstable or unlearnable target.

A search for additional, convention-compatible $\ang$ values that
could densify the $h_g(r)$ training set beyond the existing 9-angle
grid found none: every helical, box150\,nm/height300\,nm/pitch150\,nm,
physically-radius-calibrated simulation in the project already lies on
the same $\{60,65,70,72,75,80,85,87,89\}^{\circ}$ grid used above. Two
additional angles exist elsewhere in the project ($40^{\circ}$,
$50^{\circ}$) but use an incompatible, smaller box size and predate
the physical-radius calibration, so mixing them in would introduce a
box-size/radius confound; they were not used.

% ---------------------------------------------------------------
\subsection{Experiment 2 — soft-label BCE: confirmed negative result}
\label{ssec:res2}

This subsection and Section~\ref{ssec:res_oracle} address
Objective~2 (voxel-level occupancy) and the first half of
Objective~3 (root-cause diagnosis; Section~\ref{sec:intro}).

\subsubsection{Global metrics}

Training reaches a best validation IoU of $0.242$ at epoch~59
(patience exhausted at epoch~99).
BCE loss is frozen at $0.541 \pm 0.001$ throughout all 99 epochs
(Figure~\ref{fig:bce_plateau}).
The mass-conservation penalty converges rapidly
($\mathcal{L}_{\rm mass} < 0.002$ by epoch~10) but does not
lower the BCE or improve spatial discrimination.

\begin{figure}[htbp]
  \centering
  \includegraphics[width=0.85\linewidth]{pinn_bce_plateau}
  \caption{%
    The soft-label network stops learning any spatial structure
    almost immediately and instead settles for predicting the
    average film density at each angle everywhere in the volume.
    Training and validation BCE loss for the soft-label
    \film{}-conditioned \nif{} over 99 epochs.
    Loss is frozen at $0.541$ throughout, indicating that
    the network converged to the alpha-conditional mean
    $\hat{\rho} \approx \bar{\rho}(\ang)$ and cannot
    escape this equilibrium.
    The mass-conservation residual (right axis, dashed)
    converges within 10 epochs but does not perturb the BCE.
  }
  \label{fig:bce_plateau}
\end{figure}

\subsubsection{Per-alpha threshold sweep}

After training, we sweep the decision threshold $\tau \in [0.05, 0.95]$
independently per deposition angle and compare the best-$\tau$ IoU
to the trivial predictor (predict all-occupied or all-empty at
threshold $= \bar{y}_\alpha$):

\begin{table}[htbp]
  \centering
  \caption{%
    At every deposition angle, the soft-label model's best-case
    voxel-occupancy prediction is no better in practice than
    simply guessing the angle's average occupied fraction
    everywhere ($\Delta\iou < 0.003$ throughout) -- it carries no
    usable spatial information.
    Per-$\alpha$ best-threshold IoU comparison (soft-label model).
    $\bar{y}_\alpha$: ground-truth occupied fraction.
    $\iou_{\rm trivial}$: trivial-predictor IoU at optimal threshold.
    $\iou_{\rm model}$: soft-label model at optimal per-$\alpha$ threshold.
    $\Delta$: improvement over trivial.
  }
  \label{tab:soft_iou}
  \begin{tabular}{rcccc}
    \toprule
    $\ang$ (°) & $\bar{y}_\alpha$ & $\iou_{\rm trivial}$ & $\iou_{\rm model}$ & $\Delta$ \\
    \midrule
    60 & 0.001 & 0.000 & 0.000 & 0.000 (excluded) \\
    65 & 0.629 & 0.629 & 0.630 & $+0.001$ \\
    70 & 0.499 & 0.499 & 0.499 & $+0.000$ \\
    75 & 0.312 & 0.305 & 0.312 & $+0.000$ \\
    80 & 0.272 & 0.271 & 0.272 & $+0.000$ \\
    85 & 0.177 & 0.177 & 0.177 & $+0.000$ \\
    87 & 0.148 & 0.148 & 0.150 & $+0.002$ \\
    89 & 0.030 & 0.030 & 0.032 & $+0.002$ \\
    \bottomrule
  \end{tabular}
\end{table}

Every deposition angle remains a trivial predictor ($\Delta\iou < 0.003$):
the model output is equivalent to predicting the global density
for each $\ang$, with no spatial discrimination of occupied from
empty voxels.

% ---------------------------------------------------------------
\subsection{Label-quality diagnosis: the oracle predictor}
\label{ssec:res_oracle}

To determine whether the spatial learning failure is architectural
or intrinsic to the labelling scheme, we evaluate the $(\alpha,z)$-oracle
predictor (Section~\ref{ssec:oracle}) on the test set
(Figure~\ref{fig:oracle}).
\emph{Provenance note}: an earlier audit of this manuscript searched
the project codebase for a preserved, independently re-runnable
script implementing this specific $(\alpha,z)$-oracle procedure
(distinct from an unrelated corotating-frame oracle used for a
different diagnostic in the companion manuscript, which does not
implement this $(\alpha, z_{\rm bin})$ binning and cannot be
substituted for it) and found none, so the numeric outputs in this
subsection could not be re-verified from a preserved artifact at
that time. The underlying archived dataset
(\texttt{datasets/levelA\_helical\_nif/\{train,test\}\_samples.npz},
generated 2026-06-15) was subsequently located intact, and a dedicated
script (\texttt{reproduce\_alpha\_z\_oracle.py}) was written to
independently recompute both Table~\ref{tab:label_stats} and the
$\delta\iou$ values of Figure~\ref{fig:oracle} directly from this
dataset. The reproduction matches the values reported here to within
rounding (Table~\ref{tab:label_stats} entries reproduce exactly to the
reported precision; $\delta\iou$ reproduces to $\leq 0.0075$ for all
$\ang \in \{65,\ldots,87\}^{\circ}$, consistent with the $\leq 0.01$
significance threshold used throughout). This diagnostic is therefore
independently reproduced from archived data, not merely asserted; the
original run's script itself was not the one preserved, but the data
and procedure are, and the numbers stand.

\begin{figure}[htbp]
  \centering
  \includegraphics[width=0.85\linewidth]{oracle_delta_iou}
  \caption{%
    Independently reproduced from the archived dataset (see provenance
    note, Section~\ref{ssec:res_oracle}): a predictor given
    the best possible depth information (and no lateral spatial
    reasoning at all) does not beat the trivial guess, suggesting the
    missing signal may not be something a better architecture alone
    could recover from these labels.
    $\delta\iou$ (oracle minus trivial) per deposition angle.
    The oracle uses depth information $(z)$ in addition to $\ang$;
    any positive $\delta\iou$ would indicate a learnable depth signal
    in the soft labels.
    Dashed line: $\delta = 0.01$ significance threshold.
    All main production alphas (65°--87°) fall below this threshold,
    consistent with depth adding no actionable binary boundary information
    in the max-normalised soft labels.
  }
  \label{fig:oracle}
\end{figure}

The oracle achieves $\delta\iou \leq 0.007$ for all $\ang \in \{65,...,87\}^{\circ}$.
Combined with the observation that $90$--$98\%$ of soft labels at
$\ang = 65$--$75^{\circ}$ lie in the intermediate range $(0.1, 0.9)$
(Table~\ref{tab:label_stats}, visualised directly in
Figure~\ref{fig:soft_hist}), this points unambiguously to the
labelling scheme as the bottleneck.

\begin{table}[htbp]
  \centering
  \caption{%
    The occupancy labels themselves are the problem: at most
    deposition angles, the vast majority of voxels carry an
    ambiguous, neither-occupied-nor-empty value rather than a
    clean binary answer, which is why the network above had
    nothing decisive to learn from.
    Label quality statistics per deposition angle (training split).
    $f_{\rm soft}$: fraction of samples with $\rho_{\rm vox} \in (0.1, 0.9)$.
    High $f_{\rm soft}$ indicates that the max-normalised
    bead-count histogram has smeared the binary occupancy
    signal across a wide intermediate range.
  }
  \label{tab:label_stats}
  \begin{tabular}{rccc}
    \toprule
    $\ang$ (°) & $\bar{\rho}_{\rm vox}$ & $\mathrm{Var}(\rho_{\rm vox})$ & $f_{\rm soft}$ \\
    \midrule
    65  & 0.536 & 0.028 & 98.2\% \\
    70  & 0.480 & 0.034 & 96.0\% \\
    75  & 0.384 & 0.037 & 89.6\% \\
    80  & 0.331 & 0.052 & 78.0\% \\
    85  & 0.228 & 0.059 & 50.0\% \\
    87  & 0.190 & 0.056 & 44.0\% \\
    89  & 0.108 & 0.025 & 38.2\% \\
    \bottomrule
  \end{tabular}
\end{table}

\begin{figure}[htbp]
  \centering
  \includegraphics[width=\linewidth]{soft_label_histogram}
  \caption{%
    Table~\ref{tab:label_stats}'s $f_{\rm soft}$ column is not an artefact
    of binning choice: at every angle except $60^{\circ}$, most of the
    voxel mass visibly sits in the ambiguous band rather than piling up
    at 0 or 1.
    Distribution of the soft occupancy label $\rho_{\rm vox}$ on the
    training split, per deposition angle (log-scale voxel count).
    Shaded band: the ambiguous range $\rho_{\rm vox}\in(0.1,0.9)$ used to
    define $f_{\rm soft}$.
    At $\ang=60^{\circ}$ almost all voxels sit near $\rho_{\rm vox}=0$
    (empty), giving a clean, learnable binary signal; from
    $\ang=65^{\circ}$ onward, the distribution broadens into the shaded
    band and a binary classifier has little unambiguous signal left to
    fit, exactly the effect discussed below.
  }
  \label{fig:soft_hist}
\end{figure}

The mechanism is as follows.  The BCE gradient
$\partial\mathcal{L}/\partial\hat{\rho} = (\hat{\rho} - y) / (\hat{\rho}(1-\hat{\rho}))$
is maximised when $y \in \{0,1\}$ and vanishes at $y = 0.5$;
the sigmoid output $\hat{\rho} = 0.5$ is simultaneously the
maximum-entropy solution and the zero-gradient equilibrium for
a uniform label $y = 0.5$.
When $98\%$ of labels fall in $(0.1, 0.9)$, the network correctly
minimises BCE by converging to the per-$\alpha$ label mean
$\hat{\rho} \approx \bar{\rho}(\ang)$, which yields the density
calibration result but zero spatial discrimination.
This is not a failure of optimisation or architecture capacity;
it is the globally optimal solution given the soft-label distribution.

% ---------------------------------------------------------------
\subsection{Experiment 3 — hard-label BCE: restoring spatial learning}
\label{ssec:res3}

This subsection addresses the remedy half of Objective~3: having
identified the labelling scheme as the root cause
(Section~\ref{ssec:res_oracle}), we test the most direct fix —
removing the ambiguity by binarising the target — and evaluate how
far it goes toward restoring voxel-level spatial prediction.

\subsubsection{Training dynamics}

Hard binary labels ($y \in \{0,1\}$, threshold $= 0.5$) restore
the BCE gradient to its full magnitude.
Starting from a fresh random initialisation, BCE decreases from
$0.669$ (epoch~1) to $0.592$ (best at epoch~50, patience exhausted
at epoch~90), with non-zero spatial IoU emerging from epoch~7
($\iou = 0.145$ at threshold~0.5)
(Figure~\ref{fig:hardlabel_curve}).

\begin{figure}[htbp]
  \centering
  \includegraphics[width=0.85\linewidth]{hardlabel_iou_vs_epoch}
  \caption{%
    Once given unambiguous occupied/empty labels, the network
    genuinely learns during training (loss keeps decreasing)
    rather than plateauing immediately as it did with the
    soft labels, though this does not by itself guarantee useful
    voxel-level spatial resolution (see the per-angle results below).
    Validation BCE and IoU (threshold~0.5) versus training epoch for
    the hard-label \film{}-conditioned \nif{} (300-epoch budget,
    patience~=~40, best at epoch~50).
    For reference, the soft-label BCE (frozen at 0.541) and best
    soft-label global IoU (0.242, all-trivial) are shown as horizontal
    dashed lines.
    BCE decreases continuously from $0.669$ to $0.592$, confirming
    genuine learning; IoU at fixed threshold~0.5 fluctuates between
    $0.11$--$0.18$ due to class imbalance.
  }
  \label{fig:hardlabel_curve}
\end{figure}

\subsubsection{Per-alpha threshold sweep}

After training, the per-$\alpha$ threshold sweep yields:

\begin{table}[htbp]
  \centering
  \caption{%
    With clean binary labels, non-trivial voxel-level spatial
    discrimination is recovered only for the sparsest, most
    open film ($\ang=89^{\circ}$); the denser films still yield no
    better than the trivial guess, showing hard labels are
    necessary but not yet sufficient for connectivity-relevant
    spatial prediction across the full angular range.
    Per-$\alpha$ best-threshold IoU for the hard-label model.
    $f_{\rm occ}^{\rm test}$: occupied fraction in balanced test set.
    $\iou_{\rm trivial}$: trivial baseline (predict-all-1 in test set = $f_{\rm occ}^{\rm test}$).
    $\iou_{\rm hard}$: hard-label model at optimal per-$\alpha$ threshold.
    $\Delta$: improvement over trivial.
  }
  \label{tab:hard_iou}
  \begin{tabular}{rcccc}
    \toprule
    $\ang$ (°) & $f_{\rm occ}^{\rm test}$ & $\iou_{\rm trivial}$ & $\iou_{\rm hard}$ & $\Delta$ \\
    \midrule
    65 & 0.544 & 0.544 & 0.544 & $+0.000$ \\
    70 & 0.511 & 0.511 & 0.511 & $+0.000$ \\
    75 & 0.421 & 0.421 & 0.421 & $-0.000$ \\
    80 & 0.420 & 0.420 & 0.420 & $-0.000$ \\
    85 & 0.378 & 0.378 & 0.378 & $-0.000$ \\
    87 & 0.361 & 0.361 & 0.364 & $+0.003$ \\
    89 & 0.084 & 0.084 & \textbf{0.113} & $\mathbf{+0.029}$ \\
    \bottomrule
  \end{tabular}
\end{table}

Only $\ang = 89^{\circ}$ achieves non-trivial spatial discrimination
($\Delta\iou = +0.029$, Table~\ref{tab:hard_iou}); all denser alphas
remain trivial predictors at optimal threshold.
The model's sigmoid outputs are narrowly banded per alpha
(spread $\approx 0.06$--$0.17$, measured as the $p_{90}-p_{10}$ range):
at $\ang = 65^{\circ}$--$85^{\circ}$, optimal thresholding at $\tau \leq 0.01$
reduces to predict-all-positive, yielding $\iou = f_{\rm occ}^{\rm test}$.
Thus the hard-label model converges to the same qualitative solution as
the soft-label model — alpha-conditional mean density with weak spatial
modulation — but via a different mechanism: the FiLM conditioning
captures the per-$\alpha$ class balance rather than the binary spatial
boundary.
Two confounds are identified: (i)~multi-seed label conflicts at
$\ang = 80^{\circ}$ and $85^{\circ}$, where three independent seeds contribute
conflicting binary labels at the same $(x,y,z)$ coordinate,
making the spatial target inconsistent for those alphas;
(ii)~balanced sampling in the test set inflates $f_{\rm occ}$ well
above the true film occupancy, raising the trivial bar relative
to a full-grid evaluation.

% ============================================================
\section{Discussion}
\label{sec:discussion}

\subsection{Positive full-field result within a narrow angular window, and its limits}
\label{ssec:disc_scalar}

The mean spatial Pearson~$r = 0.984$ in Experiment~1 confirms that
a \film{}-conditioned \nif{} \emph{variant} (\nif{}~v5,
Section~\ref{ssec:exp1} — a different, wider network than the one
used in Experiments~2--3) is a viable full-field density surrogate
within the narrow $[75^{\circ},87^{\circ}]$ window it was trained and evaluated
on. This is useful for fast parameter sweeps restricted to that
window, but it does not establish a result at the wider
$\ang \in \{60,\ldots,89\}^{\circ}$ range studied in Experiments~2--3,
and it does not imply that the shared Experiments~2--3 architecture
(Section~\ref{ssec:arch}) — a narrower, differently-conditioned
network trained on a disjoint dataset — inherits this result:
the two are not the same model. Experiment~1's within-window
success also does not resolve the spatial-occupancy question
addressed by Experiments~2--3: it reconstructs a full 3D density
field, but always within the same 7 trained angles, so it says
nothing about a network's ability to generalise \emph{across} the
full deposition-angle range or to a disjoint angle set.

The held-out interpolation test (Section~\ref{sssec:res1_heldout},
Table~\ref{tab:nifv5_heldout}) sharpens this limit further: it is not
only that the model is untested outside $[75^{\circ},87^{\circ}]$, but
that it does not interpolate \emph{inside} that window either. The
\film{} generator (Section~\ref{ssec:arch}) takes the deposition angle
as a continuous scalar input and, in principle, nothing prevents it
from producing a smoothly-varying gamma/beta modulation for an
untrained angle between two trained ones. In practice, at $r\approx0$
for all three interpolated angles, the network behaves as if it has
allocated a separate solution to each of the seven trained angles
rather than fitting one continuous function of $\ang$ -- effectively a
lookup table implemented in weight space, not a generalisable field.
This is not an artefact of the narrow $[75^{\circ},87^{\circ}]$ window:
Section~\ref{sssec:res1_widerange} reruns the identical recipe on an
independent dataset spanning the full $[60^{\circ},89^{\circ}]$ range
and finds the same failure mode ($r=0.088$ mean at two held-out
angles, versus $r=0.954$ in-sample), so the lookup-table behaviour
appears to be a property of this \film{}-conditioning recipe itself
rather than of the specific narrow-window dataset it was first
observed on.
The explicit angle-smoothness penalty applied during training
(finite-difference regularisation of the output at $\ang \pm 0.25^{\circ}$,
Section~\ref{ssec:arch}) was intended to discourage exactly this
failure mode and evidently did not succeed at the length scale
needed here (training angles are spaced $1.25$--$2.5^{\circ}$ apart,
an order of magnitude coarser than the $0.25^{\circ}$ smoothness
window). We tested the length-scale mismatch directly: retraining the
identical recipe with the smoothness window widened to $1.5^{\circ}$
(matched to the training-angle spacing) reaches comparable in-sample
fit quality (mean $r=0.946$ across the seven trained angles, versus
$0.984$ for the original $0.25^{\circ}$ window) but the held-out mean
$r$ remains effectively zero ($r=0.006$ across the three held-out
angles). A simple length-scale widening of this penalty is therefore
not, on its own, a sufficient fix; the lookup-table behaviour persists
even once the smoothness window is matched to the data's angular
resolution. The practical consequence for Objective~1 is that \nif{}~v5
should be read strictly as a per-angle field cache with a convenient
shared architecture, not as an angle-continuous surrogate: it is
useful for fast recall at the 7 simulated angles but cannot be relied
on to interpolate film structure at an arbitrary intermediate angle
without a simulation at that angle to train (or validate) against.

\subsection{The soft-label equilibrium trap}
\label{ssec:disc_softtrap}

The frozen BCE at 0.541 is consistent with the network converging
to the correct solution of the wrong optimisation problem, rather
than a conventional training failure.
Max-normalised bead-count histograms are designed to preserve
the full density gradient (from void to densely packed), and
this is precisely why they are poor targets for binary BCE.
The information-theoretic argument is simple: if $f_{\rm soft} \approx 0.98$,
the effective entropy of the label distribution is near $H(0.5) = 1\,\mathrm{bit}$,
and a network that outputs the per-$\alpha$ mean
$\hat{\rho} \approx \bar{\rho}(\ang)$ achieves BCE $\approx 0.541$
without any spatial computation.
The mass-conservation penalty corrects the systematic per-$\alpha$
density bias but cannot force the BCE below this equilibrium.
We stop short of calling this plateau a proof of representational
impossibility: a related plateau in a separate, later production
NIF trained on this same simulation family
(pooled multi-angle/multi-seed, stalling at the maximally
uninformative $\ln 2$ baseline) was tested directly via a
controlled single-realization overfit in the companion manuscript
(Section~\ref{sec:intro}) and found representable to
near-perfect accuracy (IoU up to 1.00) once the confounds
(early stopping, pooled contradictory multi-seed labels) were
removed — i.e.\ that plateau was \emph{not} representational
impossibility. We did not run the analogous overfit test for the
soft-label equilibrium reported here, so we present the
information-theoretic argument above as the most likely
explanation, not as an experimentally isolated one.
This is consistent with, not contradicted by, the companion
result: a single deterministic, non-pooled realization is exactly
the case in which the soft-label equilibrium trap described above
does not bind, since the pathological averaging over contradictory
multi-seed soft labels this section identifies as the bottleneck
has no opportunity to occur when train and test are the same
realization.

\subsection{The oracle as an architecture-independent test}
\label{ssec:disc_oracle}

We regard the $(\alpha,z)$-oracle as the most valuable methodological
idea in this paper's diagnostic toolkit. Its numeric outputs were
independently reproduced from the archived dataset
(provenance note, Section~\ref{ssec:res_oracle}), so the discussion
below refers to a finding confirmed on the exact numbers reported
here, not merely to what the test is designed to show in principle.

The $(\alpha,z)$-oracle serves as an upper bound on the
spatial information available to any model restricted to the same
$(\alpha, z_{\rm bin})$ resolution the oracle uses, without any
spatial $(x,y)$ discrimination.
The near-trivial $\delta\iou$ ($\leq 0.0075$ for $\ang \in
\{65,\ldots,87\}^{\circ}$, reproduced above) is evidence that a
model conditioned only on this depth/angle information cannot
exceed trivial spatial discrimination on these soft labels; it does
not by itself rule out every possible architecture or conditioning
scheme (e.g.\ one given direct $(x,y)$ structural priors the oracle
does not have access to), since the oracle is a bound on
\emph{this specific} information channel, not an exhaustive
architecture search.
An architecture achieving $\Delta\iou > 0.05$ on these soft labels
would need to exploit information outside what this oracle
captures — plausible in principle, not tested here.
This diagnostic's qualitative conclusion that the labelling scheme rather than
architecture is the bottleneck is further corroborated by the hard-label
remedy in Experiment~3 (Section~\ref{ssec:res3}), which restores
genuine spatial learning without any architecture change.

\subsection{Hard labels as a minimum prerequisite}
\label{ssec:disc_hardlabel}

Binarising the density field at 0.5 is the minimal intervention
that restores binary boundary information.  It introduces a
trade-off: the continuous density gradient (useful for effective-medium
models) is discarded in favour of a crisp occupied/empty label.
For applications that require only the binary pore skeleton
(connectivity analysis, geometric surface area), hard labels are
the natural choice.
For applications requiring a continuous density field, a
full-field NIF surrogate remains a viable option \emph{within a
trained angular window} — Experiment~1 demonstrates this for its
own $[75^{\circ},87^{\circ}]$ window with its own separately-configured network
(Section~\ref{ssec:exp1}), not for the soft-label network of
Experiment~2 itself, which was not evaluated as a density
surrogate; only spatial binary structure requires hard labels.

A direct follow-up combining the two remaining prescriptions named
above — seed-isolated splits and focal loss (the class-imbalance-focused
loss reweighting of \citet{lin2017focal}, motivated here by the same
occupied/void imbalance underlying the soft-label trap) — was
subsequently run (early-stopped at epoch~291). It did not resolve spatial learning:
per-$\alpha$ inspection showed the model remained statistically
indistinguishable from the trivial predictor at 5/7 angles
($\Delta\iou < 0.0001$ at $\ang = 65$--$85^{\circ}$), with only modest
improvement at $\ang = 87^{\circ}$ ($\Delta\iou = 0.005$) and
$\ang = 89^{\circ}$ ($\Delta\iou = 0.058$); the aggregate IoU (0.419) looks
positive only because it is dominated by these few angles and is
not representative of per-angle performance. This indicates
multi-seed label pooling is not, on its own, the dominant
confound, and a higher-capacity local spatial encoding — the one remaining untested
prescription — is the more promising direction for future work.

\subsection{Multi-seed pooling at $\ang = 80^{\circ}$, $85^{\circ}$}
\label{ssec:disc_multiseed}

Pooling three independently seeded films at the same $\ang$
creates an additional secondary confound: voxels at the same
$(x,y,z,\ang)$ coordinate can have different binary occupancy
labels across seeds, since GLAD column nucleation is
seed-dependent.
For soft labels, pooling averages these into intermediate values
(further pushing labels into the soft zone).
For hard labels, pooling creates contradictory binary targets
at the same spatial coordinate.
A follow-up run applying exactly this isolation (combined with
focal loss) was carried out and is reported in
Section~\ref{ssec:disc_hardlabel}/Conclusion: it did not resolve
spatial learning at most angles, indicating multi-seed pooling is
not the dominant confound on its own.

\subsection{Testing the remaining prescriptions: spatial encoding, loss framing, and partial-realization conditioning}
\label{ssec:disc_tier2}

The seed-isolation ablation above left a higher-capacity local spatial encoding as the one
untested prescription among the three originally proposed. We tested it,
together with two further candidate fixes, directly. All three share the same
pooled hard-label dataset and evaluation protocol (per-$\ang$ threshold-swept
IoU against the trivial predict-all-occupied baseline) as the seed-isolation
run above, so any improvement is attributable to the specific change tested,
not to a different data split.
Figure~\ref{fig:three_fixes} summarises the three remedies and the
distinct reason each one fails, before the detailed discussion below.

\begin{figure}[htbp]
  \centering
  \includegraphics[width=\linewidth]{three_fixes_summary}
  \caption{%
    Three structurally different remedies were tested, not one remedy
    tried three times: a richer encoder, a different loss/target, and a
    different task framing. That all three fail, each for a distinct and
    separately diagnosed reason, is stronger evidence of a genuine
    non-identifiability than any single failed attempt would be.
    Schematic summary of the three remedies tested in this section
    (conceptual diagram; the underlying numbers are real and detailed in
    the text and Table~\ref{tab:fixc_permtest}).
  }
  \label{fig:three_fixes}
\end{figure}

\textbf{Fix A (higher-capacity local spatial encoding): a multi-resolution hash-grid encoder}
(Instant-NGP-style multi-resolution hashing, 8 levels, table size $2^{12}$/level,
resolutions 16--256) replacing the global Fourier positional encoding, all
other architecture/training settings unchanged. We use this term deliberately
rather than ``structural inductive bias'': a multi-resolution hash grid
increases local coordinate capacity and training speed, but --- unlike a
translation-equivariant architecture such as a CNN or GNN --- it imposes no
structural prior linking predictions at nearby spatial locations across
different inputs; it is a capacity change, not an equivariance constraint.
Result: no meaningful
improvement over the pooled hard-label baseline at any angle, including the
two multi-seed angles ($\Delta\iou = 0.0001$ at both $\ang=80^{\circ}$ and $85^{\circ}$,
statistically indistinguishable from the baseline's own $\Delta\iou \approx
0.0000$ there). Richer local spatial capacity alone does not resolve the
non-identifiability.

\textbf{Fix B (loss/target reframing): continuous-density regression} (MSE
loss on the un-binarised density field, same architecture and dataset as
Experiment~2's soft-label run, threshold-binarised only for the IoU
comparison). Result: no meaningful improvement ($\Delta\iou = 0.0001$ and
$0.0003$ at $\ang=80^{\circ}$ and $85^{\circ}$ respectively), and a small regression at
$\ang=60^{\circ}$ and $89^{\circ}$ relative to the hard-label baseline. The BCE-on-continuous-
target loss mismatch identified in Section~\ref{ssec:disc_softtrap} is therefore not,
by itself, the binding constraint once labels are hard.

\textbf{Fix C (task reframing): partial-realization conditioning.} Rather than
asking the network to recover a specific realization's occupancy from
$(x,y,z,\ang)$ alone -- which Section~\ref{ssec:disc_multiseed}'s oracle
diagnostic already shows has no realization-invariant target to recover, since
column positions are an independent per-realization draw -- Fix C conditions
the prediction of a realization's upper ($z\in[23,151)$) target region on a
CNN-encoded latent of that \emph{same} realization's own near-substrate
($z\in[0,23)$) observed region, reframing the task as completing a partially
observed instance of one specific structure rather than predicting an
unrelated random draw. We validate this with a 3-seed causal-substitution
permutation test at $\ang=80^{\circ}$ and $85^{\circ}$ (the two angles with 3 independent
seeds each): for every query realization, the target-region prediction is
computed once conditioned on its own observed region (same-realization) and
once for every other same-$\ang$ seed's observed region substituted in
(cross-realization); a working causal mechanism should show significantly
higher same-realization agreement (vertical-profile Pearson correlation) than
cross-realization. \textbf{Result: Fix C fails this test at both angles}
(Mann--Whitney $p=0.36$ at $\ang=80^{\circ}$, $p=0.55$ at $\ang=85^{\circ}$ -- same- and
cross-realization correlations are statistically indistinguishable, both
$\approx 0.05$--$0.07$). Diagnosing why, rather than reporting only the
failure, reveals two distinct causes. At $\ang=80^{\circ}$, the observed region
itself carries almost no seed-identifying information: per-seed occupancy in
that region is $99.80$--$99.84\%$ (pairwise pixel difference
$<0.4\%$) -- the near-substrate slab is essentially saturated regardless of
seed by the time nucleation reaches that height, so there is nothing
meaningful to condition on; this is a limitation of the observed-region
definition at this angle, not of the model. At $\ang=85^{\circ}$, the observed region
\emph{does} carry real seed information (occupancy $91.5$--$93.2\%$, pairwise
pixel difference $13$--$15\%$, an order of magnitude more variation than at
$80^{\circ}$), yet the trained encoder still fails to exploit it: inspecting the
encoder's own latent output directly shows same- versus cross-seed downstream
predictions differ by only $\sim 10^{-3}$ in magnitude. The baseline
conditioning pathway's modulation is deliberately damped to a maximum $10\%$
perturbation (inherited, for training stability, from the $\ang$-only
baseline architecture), so we tested directly whether this damping, rather
than insufficient data, was the cause: retraining with the modulation fully
unclamped (a $10\times$ larger allowed influence on the output) leaves the
permutation test result unchanged (Mann--Whitney $p=0.27$ at $\ang=80^{\circ}$,
$p=0.73$ at $\ang=85^{\circ}$, if anything slightly worse), with same/cross
prediction differences remaining $\sim 10^{-3}$ in magnitude despite the
larger budget. This rules out conditioning-pathway capacity as the
bottleneck: with only three realizations sharing this angle, the
observed-region encoder simply does not have enough training signal to learn
a useful seed-discriminative embedding, independent of how much downstream
influence it is permitted -- a genuine data-scale limitation, not evidence
that partial-realization conditioning cannot work in principle.

\textbf{Testing the data-scale hypothesis directly.} The diagnosis above
predicts that adding realizations at $\ang=85^{\circ}$ should improve
same-versus-cross separability. We tested this: five additional
simulations were run at $\ang=85^{\circ}$ (seeds 3--7, identical protocol and bead
radius to the original three), giving $n=8$ realizations at this angle
(up from $n=3$) and re-running both Fix C variants from scratch on the
enlarged dataset. $\ang=80^{\circ}$ was not extended, since its observed-region
occupancy is already $\geq 99.8\%$ at every existing seed -- a ceiling on
the observed-region definition itself that additional seeds cannot relieve.
\textbf{Result: Fix C still fails at $n=8$}, for both the damped and
unclamped variants (Mann--Whitney $p=0.35$ and $p=0.53$ respectively, up
from $6$ to $56$ cross-realization pairs), with same- and
cross-realization correlation means remaining statistically
indistinguishable ($0.127$ vs.\ $0.126$ damped; $0.124$ vs.\ $0.122$
unclamped) -- not a near-miss, but a flat null result at substantially
higher statistical power. More than doubling ($n=3\to8$) the per-angle
realizations therefore does \emph{not} resolve Fix C's failure at
$\ang=85^{\circ}$: the data-scarcity explanation, while a reasonable hypothesis
from the diagnostics above, is not supported once tested directly at the
scale practically available to this project. The bottleneck -- real
seed-specific signal is demonstrably present in the observed input
($9$--$15\%$ pixel-level variation across seeds) yet is not exploited by
the trained encoder at any tested sample size or FiLM-damping setting --
should be reported as a genuinely unresolved architecture/objective
limitation rather than attributed to a data-volume shortfall that a modest
seed increase was expected to fix. We report this as a strengthened
negative result: it directly pre-empts, rather than merely anticipates,
the natural reviewer objection that more seeds would resolve the
non-identifiability. (The 56 cross-realization pairs are clustered --
eight query-realizations with seven sources each, and predictions vary
only weakly across sources for a fixed query, which is itself the
observed inertness of the conditioning pathway; this clustering, if
anything, biases the test toward detecting a difference, so the null
result is conservative.)
Table~\ref{tab:fixc_permtest} summarises all six causal-substitution
permutation tests run for Fix C across both FiLM-damping variants,
both tested angles, and both sample sizes: none reaches significance,
and the $n=3 \to n=8$ extension at $\ang=85^{\circ}$ leaves the same- and
cross-realization correlation means essentially unchanged while
increasing the cross-realization sample by nearly an order of
magnitude.

\begin{table}[htbp]
  \centering
  \caption{%
    Six independent tests, two FiLM-conditioning strengths, two sample
    sizes -- and not one shows a same- versus cross-realization
    separation a working causal mechanism would produce.
    Fix C causal-substitution permutation test results (Mann--Whitney
    $U$ on same- versus cross-realization vertical-profile Pearson
    correlation). $n$: independent seeds at that angle; ``pairs'':
    same-realization / cross-realization comparison counts.
    Unclamped $n{=}3$ same-/cross-realization means were not logged
    separately (only individual example values and the $p$-value were
    recorded); the text (Section~\ref{ssec:disc_tier2}) reports both
    populations in the $\approx 0.05$--$0.07$ range, consistent with
    the damped $n=3$ values shown here.
  }
  \label{tab:fixc_permtest}
  \begin{tabular}{llccccc}
    \toprule
    Variant & $\ang$ (°) & $n$ & pairs (same/cross) & same-$r$ & cross-$r$ & $p$ \\
    \midrule
    Damped     & 80 & 3 & 3/6  & 0.073 & 0.073 & 0.36 \\
    Damped     & 85 & 3 & 3/6  & 0.055 & 0.055 & 0.55 \\
    Unclamped  & 80 & 3 & 3/6  & \multicolumn{2}{c}{$\approx$0.05--0.07} & 0.27 \\
    Unclamped  & 85 & 3 & 3/6  & \multicolumn{2}{c}{$\approx$0.05--0.07} & 0.73 \\
    Damped     & 85 & 8 & 8/56 & 0.127 & 0.126 & 0.35 \\
    Unclamped  & 85 & 8 & 8/56 & 0.124 & 0.122 & 0.53 \\
    \bottomrule
  \end{tabular}
\end{table}

\subsection{Architecture-independent cross-validation: classical reconstruction and global connectivity agree with the diagnosis}
\label{ssec:disc_crossval}

Every remedy tested above (Fix A/B/C, and the seed-isolation ablation) is a
point-wise, coordinate-conditioned neural model. A reasonable objection is
that the failure is specific to that model family -- some other, non-neural
approach might recover the missing structure. We tested this directly with
two architecture-independent checks that use no coordinate-conditioned neural
network at all: a classical stochastic-reconstruction baseline, and an
independently-computed global connectivity descriptor from a separate
simulation campaign. Both converge on the same conclusion as Fix C.

\paragraph{A classical stochastic-reconstruction baseline.}
Two-point correlation functions do not, in general, uniquely determine a
microstructure realization~\citep{yeong1998reconstructing} -- but the
size of that gap is an empirical question for any specific material system,
not a foregone conclusion. We measured it directly for these films using the
Yeong--Torquato swap-based simulated-annealing algorithm: for each
$\ang$, every $z$-slice of the true occupancy volume is reconstructed from
scratch, starting from a random binary field at the correct volume fraction
$\phi$ and repeatedly swapping an occupied/unoccupied voxel pair whenever the
swap reduces the squared error between the candidate's radially-averaged
two-point correlation $S_2(r)$ (periodic boundary, $r=1$--$12$ voxels, short
range $r\!\leq\!3$ weighted $4\times$) and the true slice's $S_2(r)$, at fixed
temperature schedule ($T_0=0.004$, geometric cooling, $18{,}000$ swaps per
slice). Because every accepted move is a $1\leftrightarrow0$ swap, $\phi$ is
matched exactly by construction; only $S_2(r)$ is optimized. Reconstructed
slices are then scored against the true slice with the same IoU metric used
throughout Table~\ref{tab:hard_iou}.

\textbf{Scope caveat, stated explicitly:} this is a different task from the
one Fix A/B/C were scored on. Fix A/B/C fit one coordinate function
$f(x,y,z,\ang)\!\to\!$occupancy jointly across all nine $\ang$, evaluated on
held-out voxel \emph{samples} pooled across angles. The reconstruction here
instead asks a complementary, arguably more information-scarce question: with
\emph{no} cross-angle training and \emph{no} coordinate signal at all -- only
the volume fraction and two-point correlation of one realization, measured
on that same realization -- how much of it can be recovered? It is an
information-theoretic ceiling test, not a drop-in fourth entry in
Table~\ref{tab:hard_iou}.

Table~\ref{tab:yt_baseline} gives the result at all nine angles ($25$ of
$151$ $z$-slices sampled per angle, seed~0). At every non-degenerate angle,
the swap-based search closes $70$--$96\%$ of the initial random-field
$S_2$-matching error -- i.e., it finds a configuration whose two-point
correlation function is very close to the true one -- yet IoU against the
true realization remains low and falls monotonically as $\phi$ decreases.
$\ang=60^{\circ}$ is a degenerate near-empty case ($\phi\!\approx\!0.002$)
included for completeness. \textbf{A classical algorithm operating directly
on nearly-exact second-order statistics, with no architecture to blame,
cannot recover the realization either.} This rules out a specific class of
counter-explanation for the neural failure (that the tested architectures
were simply the wrong ones) without weakening the diagnosis: the missing
information is not resident in any second-order statistic of the film, for
any method tested.

\begin{table}[htbp]
  \centering
  \caption{%
    Yeong--Torquato classical stochastic-reconstruction baseline, all nine
    trained angles (seed~0, $25/151$ $z$-slices sampled per angle).
    $S_2$ error closure: $1 - (\text{final squared } S_2 \text{ error} /
    \text{initial random-field squared } S_2 \text{ error})$, i.e.\ how much
    of the gap to a perfect two-point-correlation match was closed by the
    swap-based search. IoU computed against ground truth on the same
    sampled slices. $\ang=60^{\circ}$ is near-empty
    ($\phi\!\approx\!0.002$) and included only for completeness.
  }
  \label{tab:yt_baseline}
  \begin{tabular}{lccc}
    \toprule
    $\ang$ (°) & $\phi$ (occupancy) & $S_2$ error closure & IoU \\
    \midrule
    60 & 0.000 & --  (degenerate) & 1.000 \\
    65 & 0.612 & 93\% & 0.468 \\
    70 & 0.486 & 95\% & 0.338 \\
    72 & 0.458 & 96\% & 0.317 \\
    75 & 0.278 & 94\% & 0.172 \\
    80 & 0.221 & 92\% & 0.132 \\
    85 & 0.121 & 85\% & 0.068 \\
    87 & 0.072 & 70\% & 0.046 \\
    89 & 0.011 & -- (already near target) & 0.007 \\
    \bottomrule
  \end{tabular}
\end{table}

\paragraph{Cross-referencing an independently-computed global connectivity descriptor.}
A separate simulation campaign (27 runs: all nine $\ang$ $\times$ three
seeds, frozen periodic connected-component analysis, $0.5$\,nm voxels,
film band $z\in[10\text{nm}, 0.90H]$) computed the largest connected-component
fraction (LF) -- the coarsest possible topology descriptor -- independently
of both P2's NIF work and the reconstruction above. Table~\ref{tab:topology_crossref}
reports, for each $\ang$, the seed-to-seed spread of LF and of volume
occupancy from that campaign, alongside whether Fix C tested that angle
directly.

\begin{table}[htbp]
  \centering
  \caption{%
    Global connectivity (largest-fraction, LF) and occupancy, mean and
    seed-to-seed spread (max$-$min over 3 seeds), from an independently-run
    27-run topology sweep. LF spread is $0.00000$ up to $\ang=80^{\circ}$ and
    reaches only $0.013$ at the most stochastic angle tested
    ($89^{\circ}$) -- versus the $9$--$15\%$ \emph{pixel-level} seed
    disagreement Fix C measured directly on the observed region at
    $\ang=85^{\circ}$ (Section~\ref{ssec:disc_tier2}).
  }
  \label{tab:topology_crossref}
  \begin{tabular}{lcccc}
    \toprule
    $\ang$ (°) & occupancy (mean, spread) & LF (mean, spread) & Fix C tested? \\
    \midrule
    60--75 & 0.63--0.83, $\leq$0.003 & 1.0000, 0.00000 & no \\
    80     & 0.500, 0.0024           & 1.0000, 0.00000 & yes -- saturated, no signal \\
    85     & 0.304, 0.0032           & 0.9999, 0.00015 & yes -- real pixel-level signal, unexploited \\
    87     & 0.198, 0.0018           & 0.9969, 0.00262 & no \\
    89     & 0.086, 0.0062           & 0.9822, 0.01306 & no \\
    \bottomrule
  \end{tabular}
\end{table}

Global connectivity is essentially seed-invariant at every angle tested --
literally deterministic to five decimal places up to $\ang=80^{\circ}$, and
still within $1.3$ percentage points at $89^{\circ}$ -- a roughly
$60$--$100\times$ gap relative to the $9$--$15\%$ pixel-level seed
disagreement Fix C independently measured at $\ang=85^{\circ}$.

A fourth check probes the same question one architecture class earlier and
more cheaply than a full training run: does giving a model explicit spatial
structure -- rather than richer coordinate encoding, a different loss, or
partial-realization conditioning with an MLP head -- change the outcome? A
small 3D U-Net (3 levels, 16/32/64 channels) was built to take the
coordinate-conditioned inputs used throughout this paper (voxel height and
$\ang$, broadcast constant over $x,y$ within each height, so that no ground-truth
occupancy is ever fed to the model) and predict full-volume occupancy patches.
Before any training-set evaluation, this architecture must pass a standard
sanity gate: it should be able to drive the loss to $\approx 0$ by memorizing
a single fixed $32^3$ patch. It could not. Three loss variants -- focal loss
at the paper's default weighting, focal loss re-weighted for the patch's
$13\%$ positive-voxel fraction, and plain binary cross-entropy -- all
plateaued at the same IoU ($0.47$--$0.56$), independent of the loss function
used. An ablation isolates the cause: adding two explicit $x,y$ coordinate
channels (not a viable input for real training, since it would let the model
memorize absolute position rather than generalize across patches) raises the
memorization IoU to $1.000$ in the same protocol. The mechanism is a direct
consequence of the input design, not an optimization or capacity failure:
height- and angle-only conditioning is constant over $x,y$ at fixed height, and
a translation-equivariant convolutional network maps a spatially-constant
input field to a spatially-constant output field (up to boundary effects),
so it cannot represent an $x,y$-varying occupancy texture from this input by
construction. Spatial inductive bias alone therefore does not remedy the
non-identifiability documented above -- a CNN is subject to the same ceiling
as the point-wise architectures of Table~\ref{tab:hard_iou} whenever it is
given the same coordinate-only information, and giving it more (absolute
position) would trade the diagnosis for label leakage rather than resolve it.

A related but distinct question is whether Fix C's own encoder -- which
\emph{does} have access to real, per-realization structural information (the
observed near-substrate region) rather than coordinates alone -- was limited
by its specific implementation rather than by the underlying task. Inspecting
that encoder shows it ends in a global-average-pooling layer, collapsing the
entire observed-region footprint into one vector before the FiLM-modulated
decoder ever sees it -- a bottleneck that discards exactly the per-column
spatial information the causal-substitution diagnosis in
Section~\ref{ssec:disc_tier2} showed is present at $\ang=85^{\circ}$
($9$--$15\%$ pixel-level variation across seeds). We therefore built a second
variant that avoids this specific bottleneck: a resolution-preserving 2D
encoder over the observed map, with each query point's prediction
pixel-aligned to its own $(x,y)$ location in that feature map via bilinear
sampling (local, not global, conditioning), replacing FiLM modulation with
direct concatenation. Before any full-scale training, this architecture was
required to pass the same class of sanity check as the coordinate-only CNN
above -- memorize a small subset of one realization's own points -- and it
did, reaching perfect accuracy in under $150$ steps, confirming the pathway
can carry and exploit local information when a target is directly
learnable from it. Trained on the full partial-conditioning dataset and
evaluated with the identical causal-substitution permutation test used for
the original Fix C, the result is unchanged: same- and cross-realization
vertical-profile correlations remain statistically indistinguishable at both
tested angles (Mann--Whitney $p=0.64$ at $\ang=80^{\circ}$, $p=0.49$ at
$\ang=85^{\circ}$ -- if anything weaker separation than the original
FiLM-based Fix C's $p=0.36$/$p=0.55$), with same-realization predictions
differing from cross-realization substitutions by no more than ordinary
run-to-run noise (e.g.\ $0.339$ same vs.\ $0.328$--$0.340$ across all seven
substituted sources, for a representative $\ang=85^{\circ}$ query). This
rules out the global-pooling bottleneck as \emph{the} explanation for Fix
C's failure: even a properly spatially-resolved, sanity-checked conditioning
pathway is not exploited once trained at scale, strengthening rather than
weakening the case that the observed-to-target correlation this task
requires is either too weak to detect at the realization counts available or
absent altogether in this class of stochastically-nucleated film, independent
of which of the encoder architectures tested here is used to try to extract
it.

\paragraph{Synthesis.} Five independent methods, on five independent analyses
of these films -- Fix C's causal-substitution permutation test (both its
original FiLM-based encoder and the pixel-aligned local-conditioning variant
above), the 27-run topology sweep, the Yeong--Torquato reconstruction above,
and the coordinate-ablation check on a spatially-aware architecture --
converge on one conclusion: seed-to-seed stochasticity in these films is real,
substantial at the pixel level, and essentially invisible to every low-order
descriptor tested (connectivity class, volume fraction, two-point
correlation), regardless of whether that descriptor is consumed by a
point-wise neural network, a spatially-convolutional one, a classical
optimizer, or read directly off an independently-computed statistic. The
negative result in Table~\ref{tab:hard_iou} is therefore not attributable to
the specific architectures tried in this paper; it reflects where the
recoverable information genuinely is (and is not) in this class of
stochastically-nucleated GLAD film.

\subsection{Limitations}
\label{ssec:disc_limits}

Beyond the label-quality and multi-seed confounds already
discussed above, the following scope limitations should be kept
in mind when reading the results of this paper alongside other
project documents.

\begin{enumerate}
  \item \textit{Single film height and pitch, and pitch is partly
    confounded with angle.} All results are for $h = \SI{300}{nm}$,
    pitch $= \SI{150}{nm}$; generalization to other pitches or heights is
    not tested. A targeted check projecting an $\ang=85^{\circ}$
    pitch sweep ($9$--$\SI{260}{nm}$) onto the PCA plane of the
    angle-sweep two-point descriptor $h_g(r;\ang)$ finds that varying
    pitch moves the descriptor $0.75\times$ as far along the dominant
    angle direction as the full $\ang=60$--$89^{\circ}$ sweep itself,
    with more than half the tested pitches landing within a
    seed-to-seed scatter radius of a \emph{different} angle's
    descriptor (e.g.\ the $\ang=85^{\circ}$, pitch$=\SI{9}{nm}$
    descriptor is closer to $\ang=87^{\circ}$, pitch$=\SI{150}{nm}$
    than to any other $\ang=85^{\circ}$ realization). Angle and pitch
    are therefore \emph{partly confounded} in this descriptor: the
    companion inversion result \citep{companionP1} recovers $\ang$
    correctly only under the fixed-pitch assumption of the primary
    campaign, and a joint $(\ang,\text{pitch})$ inversion would need a
    second, pitch-sensitive observable to break the degeneracy.
  \item \textit{Ballistic model only.} The density fields are from
    a diffusion-free bead simulator; real Cu GLAD films at
    $\ang > 80^{\circ}$ may have structural differences.
  \item \textit{Multi-seed label collision at $\ang = 80^{\circ}$, $85^{\circ}$.}
    Discussed in Section~\ref{ssec:disc_multiseed}.
  \item \textit{Hard-label threshold.} The 0.5 threshold for
    binarisation is physically motivated (majority occupation) but
    arbitrary; sensitivity to threshold choice is not studied here.
  \item \textit{Occupied-fraction values are $z$-range-dependent and
    not directly comparable across this project's own documents.}
    Table~\ref{tab:soft_iou}'s $\bar{y}_\alpha$ column uses the full
    $z\in[0,151)$ voxel range. A separate, later analysis restricted
    to training slices $z\in[0,121]$ reports different occupied
    fractions for the nominally same configurations (e.g.\ at
    $\ang=85^{\circ}$: $0.126$ vs.\ $0.177$ here; at $\ang=89^{\circ}$: $0.012$
    vs.\ $0.030$ here) — a $z$-range definitional difference, not a
    data inconsistency, but one that was not disclosed here and
    should not be conflated across documents without noting which
    $z$-range each value uses.
  \item \textit{Computational cost of the underlying simulations.}
    Generating each training configuration requires a dedicated
    GPU-accelerated ballistic run; the number of angle/pitch/height
    combinations that can be simulated, and hence the density and
    coverage of the NIF training set, is bounded by this cost rather
    than by any modelling choice; any future extension of the
    parameter space inherits this same constraint.
\end{enumerate}

\subsection{Outlook: Fast morphology-to-optics workflow via NIF surrogates}
\label{ssec:disc_optical_design}

The companion paper \citep{companionP1} characterises the same ballistic simulator's optical properties via finite-difference time-domain (FDTD) electromagnetic modelling of Cu GLAD films across a pitch range of $9$--$260\,\mathrm{nm}$ at $\ang=85^{\circ}$, finding that the optical response (reflectance, transmittance, absorptance across $300$--$1500\,\mathrm{nm}$) is pitch-invariant—a consequence of the columnar independence encoded in the ballistic model.
This optical validation, combined with the density-field surrogate developed here, opens a pathway to fast closed-loop optical design workflows: experimentalists could evaluate film morphology predictions from a trained NIF in milliseconds, combine those predictions with pre-computed optical lookup tables (or trained optical surrogates), and iterate on angle/pitch/height without running full FDTD simulations or experimental depositions.
Such a workflow would be particularly valuable for applications like metasurfaces, thermal absorbers, or angle-tunable optical elements, where the morphology-to-optics mapping is critical and parameter-space screening is currently prohibitive.
The present paper's demonstration that compact descriptors (pair-correlation, density profile) generalise where full fields do not suggests that analogous descriptor-based optical surrogates may also be learnable; systematic exploration of this cross-domain surrogate coupling is reserved for future work.

% ============================================================
\section{Conclusion}
\label{sec:conclusion}

We have investigated the use of a \film{}-conditioned neural implicit
field for scalar porosity and binary voxel occupancy prediction in
ballistic helical GLAD simulations, with the aim of providing an
experimentalist with a fast surrogate for process-window screening.
We close each of the three objectives stated in
Section~\ref{sec:intro}, in the same order.

\begin{itemize}
  \item \textbf{Objective 1 (scalar porosity surrogate) — answered:
    full-field density reconstruction within $[75^{\circ},87^{\circ}]$ is
    positive, but not a wider-domain result.}
    A separately-configured \nif{}~v5 variant achieves mean spatial
    Pearson $r = 0.984$ reconstructing the 3D density field at its
    7 trained angles $\ang \in \{75, 77.5, 80, 82.5, 85, 86, 87\}^{\circ}$.
    This confirms the architecture family's suitability as a
    density-level surrogate \emph{within that window}; it does not
    establish, and this experiment does not test, generalisation to
    the wider $\{60,\ldots,89\}^{\circ}$ range used in Experiments~2--3,
    nor is it the same network as the shared architecture used
    there. A held-out test at three angles inside the trained window
    but withheld from training ($78.75^{\circ}$, $83.5^{\circ}$,
    $84.25^{\circ}$; Section~\ref{sssec:res1_heldout}) finds this
    correlation collapses to $r = 0.003$ mean, showing the $r=0.984$
    figure reflects per-angle recall rather than a genuinely
    angle-continuous learned field -- the model does not interpolate
    even within its own trained window. An independent replication
    of the same recipe on the full $[60^{\circ},89^{\circ}]$ range
    (Section~\ref{sssec:res1_widerange}, disjoint dataset, 7 trained /
    2 held-out angles) confirms this is not specific to the narrow
    window: in-sample $r = 0.954$ collapses to $r = 0.088$ at the two
    held-out angles.
    Two follow-up tests (Section~\ref{sssec:res1_descriptors}) localise
    this failure to the \emph{target}, not the conditioning
    architecture. Replacing \film{} with direct Fourier-concatenation
    of $\ang$, same full-field target and split, does not fix
    generalisation ($r=0.101$ held-out, vs.\ \film{}'s $r=0.088$ --
    both effectively zero, not a meaningful difference). Reformulating
    the target instead, from the raw field to the compact
    pair-correlation descriptor $h_g(r;\ang)$ already validated by P1,
    does generalise on this wide range (mean $R^2=0.962$, Pearson
    $r=0.997$ across two held-out angles) -- though, matching P1's own
    caveat, not clearly better than naive interpolation at these same
    points. The same reframing is not a universal fix: scalar
    porosity($\ang$) is better served by classical linear interpolation
    (mean abs.\ error $0.035$) than by any network at this data volume
    (NIF: $0.181$), and the depth-density-profile descriptor
    $\rho_z(z;\ang)$'s generalisation is angle-dependent, established
    after fixing the seeding gap and averaging ten properly-seeded
    reruns: its \emph{shape} is reliably recovered ($r=0.766\pm0.041$,
    positive in all ten runs) but its absolute level only generalises
    well inside the training range ($R^2=0.457\pm0.236$ at
    $72^{\circ}$) and fails near the training range's edge
    ($R^2=-0.419\pm0.385$ at $87^{\circ}$) -- a genuine partial
    success, not an unresolved question.

  \item \textbf{Objective 2 (voxel-level occupancy) — answered:
    soft-label spatial occupancy is negative.}
    BCE with max-normalised density labels freezes at loss $= 0.541$;
    all deposition angles remain trivial predictors, i.e.\ the
    connectivity/percolation-relevant voxel-level information
    motivating Objective~2 is not recovered by this labelling
    scheme.
    The physics-informed mass-conservation penalty corrects the
    per-$\alpha$ density bias but cannot restore binary spatial
    discrimination from soft labels: it is a density-calibration
    tool only, not a fix for the spatial-occupancy problem.

  \item \textbf{Objective 3 (root cause and remedy) — answered:
    diagnosed root cause; three candidate remedies tested, none
    sufficient at the current data scale, with the reasons for each
    failure diagnosed rather than left open.}
    Root cause: $90$--$98\%$ of soft labels lie in $(0.1,0.9)$,
    eliminating binary boundary gradients from the BCE gradient;
    an $(\alpha,z)$-oracle diagnostic rules out architecture
    capacity and depth-signal absence as alternative explanations.
    Retraining with hard binary labels ($\{0,1\}$, threshold $0.5$)
    reduces BCE from $0.669$ to $0.592$ and achieves non-trivial
    IoU for $\ang = 89^{\circ}$ ($\Delta\iou = +0.029$); $6/7$ other angles
    remain trivial due to two identified confounds: multi-seed
    label conflicts at $\ang = 80^{\circ}$ and $85^{\circ}$, and
    balanced-sampling distortion of the test class distribution.
    Hard binary labels are therefore a \emph{necessary but not
    sufficient} condition.
    A direct follow-up combining seed-isolated splits and focal
    loss (early-stopped at epoch~291) was subsequently run and did
    \emph{not} resolve spatial learning: 5/7 angles remained
    statistically indistinguishable from the trivial predictor
    ($\Delta\iou < 0.0001$ at $\ang = 65$--$85^{\circ}$), with only modest
    improvement at $\ang = 87^{\circ}$ ($\Delta\iou = 0.005$) and
    $\ang = 89^{\circ}$ ($\Delta\iou = 0.058$); the aggregate IoU (0.419)
    is not representative of per-angle performance.
    This rules out multi-seed pooling as the dominant confound on
    its own. Three further remedies were then tested directly
    (Section~\ref{ssec:disc_tier2}): a multi-resolution hash-grid
    spatial encoder (higher-capacity local encoding, not a structural
    inductive bias), continuous-density
    regression (a loss/target reframing), and partial-realization
    conditioning (a task reframing, conditioning prediction of a
    realization's upper region on an encoded latent of that same
    realization's own near-substrate region). The first two produce
    no meaningful improvement over the hard-label baseline at any
    angle. The third fails a 3-seed causal-substitution permutation
    test at both multi-seed angles, for two mechanistically
    distinct reasons: at $\ang=80^{\circ}$ the near-substrate observed
    region is itself $>99.8\%$ saturated regardless of seed, leaving
    nothing informative to condition on; at $\ang=85^{\circ}$ real
    seed-distinguishing signal exists in the observed region
    ($13$--$15\%$ pairwise pixel variation) but the trained encoder's
    influence on the output remains $\sim10^{-3}$ in magnitude. A
    direct follow-up retrained the conditioning pathway fully
    unclamped ($10\times$ larger allowed influence, removing the
    $10\%$ FiLM damping inherited from the baseline architecture) to
    test whether the damping itself was responsible: the permutation
    test still fails at both angles with the same order-of-magnitude
    same/cross difference, ruling out conditioning-pathway capacity
    as the bottleneck and narrowing it, provisionally, to the number
    of realizations sharing each angle (three). This data-scarcity
    hypothesis was then tested directly: five additional simulations
    at $\ang=85^{\circ}$ (seeds~3--7) brought that angle to $n=8$
    realizations, and both Fix C variants were retrained from
    scratch on the enlarged dataset. \textbf{The permutation test
    still fails at $n=8$} for both variants (Mann--Whitney $p=0.35$
    damped, $p=0.53$ unclamped, evaluated over $56$ cross-realization
    pairs versus $6$ previously), with same- and cross-realization
    correlation means remaining statistically indistinguishable
    ($0.127$ vs.\ $0.126$; $0.124$ vs.\ $0.122$) rather than merely
    falling short of significance. More than doubling the per-angle
    realizations therefore does \emph{not} resolve Fix C's failure.
    A higher-capacity local spatial encoding, loss/target reframing, and now a
    practically-achievable increase in per-angle realizations are
    all ruled out as sufficient fixes; partial-realization
    conditioning's failure at $\ang=85^{\circ}$ -- where real seed-specific
    signal is demonstrably present in the input yet is not exploited
    by the trained encoder at any tested sample size or
    FiLM-damping setting -- is reported as a genuine, tested negative
    result rather than an artefact of insufficient data, pending
    architectural or objective-function changes not explored here.
\end{itemize}

This diagnosis is further strengthened by four architecture-independent
cross-checks (Section~\ref{ssec:disc_crossval}) that use no
coordinate-conditioned neural network at all: a classical Yeong--Torquato
stochastic-reconstruction baseline, operating directly on the two-point
correlation function of each realization, closes $70$--$96\%$ of the
initial random-field correlation-matching error yet still recovers only
$7$--$47\%$ intersection-over-union with the true realization at
non-degenerate angles; an independently-run 27-configuration topology
sweep shows the global largest-connected-fraction descriptor is
seed-invariant to five decimal places up to $\ang=80^{\circ}$ and within
$1.3$ percentage points at $\ang=89^{\circ}$, a $60$--$100\times$ gap
relative to the $9$--$15\%$ pixel-level seed disagreement measured
directly; a coordinate-conditioned convolutional network is shown, by
construction, to be representationally incapable of producing an
$x,y$-varying occupancy texture from height/angle-only input; and a
spatially-resolved, sanity-checked local-conditioning variant of the
partial-realization encoder -- built specifically to rule out the
global-average-pooling bottleneck of the original encoder -- still shows
no learned signal once trained at scale. Five independent methods, on
five independent analyses of these films, therefore converge on one
conclusion: seed-to-seed stochasticity in these films is real,
substantial at the pixel level, and essentially invisible to every
low-order descriptor tested, regardless of whether that descriptor is
consumed by a point-wise neural network, a spatially-convolutional one, a
classical optimizer, or read directly off an independently-computed
statistic. The negative result is therefore not attributable to the
specific architectures tried in this paper; it reflects where the
recoverable information genuinely is (and is not) in this class of
stochastically-nucleated GLAD film.

These results establish a stepwise label-quality diagnosis
framework — soft-zone fraction analysis, oracle predictor test,
and hard-label ablation — applicable to any simulation-based
\nif{} training pipeline where continuous density fields are the
natural simulation output but binary voxel structure is the
prediction target. For an experimentalist, the practical takeaway
is that a fast scalar-porosity surrogate is already usable within
its trained angular window, while a voxel-resolved,
connectivity-aware surrogate is not yet available: spatial inductive
encoding alone does not supply it (Section~\ref{ssec:disc_tier2}), and
the most promising remaining direction, partial-realization
conditioning, requires more per-angle realizations and a stronger
conditioning pathway than were available in this study before it can
be properly evaluated.

% ============================================================
\section*{Author contributions (CRediT)}
O.B.H.A.: Conceptualization, Methodology, Software, Investigation, Formal analysis,
Writing -- original draft. B.G.: Supervision (Paris-based work), Writing -- review \&
editing. F.C.A.: Supervision (principal), Writing -- review \& editing.

% ============================================================
\section*{Acknowledgements}
The authors thank their home institutions for their support:
the Laboratoire de Photovolta\"{i}que et Mat\'{e}riaux
Semi-conducteurs (LPMS), \'{E}cole Nationale d'Ing\'{e}nieurs de
Tunis (ENIT), Universit\'{e} de Tunis El Manar (UTM); the Institut
des NanoSciences de Paris (INSP), CNRS, Sorbonne Universit\'{e};
and the FEMTO-ST Institute, CNRS, Universit\'{e} de
Franche-Comt\'{e}.

\section*{Funding}
No specific financial support statement is declared for this manuscript
because no verified project number, contract number, or funding-agency
wording was found in the local project records.

\section*{Data availability}
Simulation data, voxel files, and training scripts are available
from the authors upon reasonable request.

% ============================================================
\bibliographystyle{aipnum4-1}
\bibliography{P2_references}

\end{document}
